\section{Sections and points: local packets, abelian fibres, and descent}
\label{sec:B}

Throughout this section $U=X\setminus S$ is a smooth geometrically connected
curve, $X$ is its smooth proper completion, $\Delta_U=\pi_1(U_{\bar k})$,
and $\Pi_U=\pi_1(U)$.  A section is continuous and is taken modulo
conjugation by $\Delta_U$.  We write
\[
 \widehat{k^\times}=\varprojlim_{n\mid m}k^\times/(k^\times)^n.
\]
This is the power completion.  Its finite-level quotients need not be finite
over a number field.  Over a finite extension of $\mathbb Q_p$ they are
finite.  These two assertions use the same inverse system and do not
identify its topology with an unspecified completion.

\subsection{Injectivity, including the point germs of cote 143}

\begin{statement}{B.1}{Proved}{The corrected semi-abelian Kummer input}
Let $k$ be finitely generated over $\mathbb Q$ and let
$1\to T\to G\to A\to1$ be a semi-abelian variety over $k$.
Then
\[
 \bigcap_{n\geq1}nG(k)=0,
 \qquad G(k)\longrightarrow H^1_{\rm cts}(G_k,TG)
 \quad\hbox{is injective}.
\]
Consequently, for an immersion $U\hookrightarrow G$, two $k$-points with
the same abelianized section class are equal.  The corresponding map on
algebraic points and germs of sections is injective.
\end{statement}
\begin{proof}
First, $\bigcap_n(k^\times)^n=\{1\}$.  Let $F$ be the algebraic closure of
$\mathbb Q$ in $k$, a number field.  On a normal proper integral $F$-model
of $k$, the valuation of an element admitting every root is divisible by
every positive integer at every prime divisor, hence is zero.  A rational
function with zero divisor on this model is a nonzero global regular
function and therefore belongs to $F^\times$.  All its roots in $k$ then
belong to $F$, by the definition of $F$.  Valuations at finite primes of
$F$ reduce the assertion to the unit group of $F$, which is a finite group
plus a finitely generated free abelian group.  The intersection of its
$n$th powers is $1$.  A finite splitting field for $T$ therefore proves
$\bigcap_n nT(k)=0$.

The Lang--N\'eron theorem and the Mordell--Weil theorem for the constant
trace give finite generation of $A(k)$.  Let $a\in\bigcap_n nG(k)$.
Its image in $A(k)$ is zero.  Let $M$ annihilate $A(k)_{\rm tors}$.
For each $n$, choose $b_n$ with $nb_n=a$.  Then $Mb_n\in T(k)$ and
$Ma=n(Mb_n)$.  Thus $Ma=0$.  The torsion subgroup of $G(k)$ is finite:
its image is contained in the finite group $A(k)_{\rm tors}$ and its
kernel in $T(k)_{\rm tors}$, which is finite after a finite splitting
extension.  Since every $b_n$ is now torsion, taking $n$ equal to the
exponent of $G(k)_{\rm tors}$ gives $a=0$.

The multiplication-by-$n$ Kummer sequence sends $a$ to zero precisely
when $a\in nG(k)$.  Passing over all $n$ proves the asserted injectivity.
The continuous-cohomology comparison is legitimate: the groups of
$k$-rational $n$-torsion are finite, so the degree-zero inverse system
is Mittag--Leffler and contributes no $\varprojlim^1$ term.
For $x,y\in U(k)$, equality of their abelian section classes gives zero
Kummer class for their difference in $G(k)$, hence $x=y$.
Finally, equality of two section germs means equality after restriction
to $G_L$ for one finite extension $L/k$.  Apply the same argument over
$L$.  This proves exactly the germ assertion in cote 143, scans 4--5,
including its stated semi-abelian embedding hypothesis.
\end{proof}


\begin{statement}{B.2}{Refuted}{Excluded types and explicit failures over larger fields}
For $\mathbb P^1$ and $\mathbb A^1$ in characteristic zero the geometric
fundamental group is trivial, so every rational point gives the same
section.  On the hyperbolic curve
$U_0=\mathbb P^1\setminus\{0,1,\infty\}$ over $\mathbb R$, the distinct
points $2$ and $3$ give the same section class.  For every prime $p$, put
\[
 k=\mathbb Q_p((t^{\mathbb Q})),\qquad x=2,\qquad y=2+t.
\]
This is a $p$-adically closed field which is not sub-$p$-adic, and $x,y$
give the same full arithmetic section class on $U_0/k$.
\end{statement}
\begin{proof}
The first assertion follows from the geometric simply connectedness of
the projective and affine line in characteristic zero.  The real path
from $2$ to $3$ inside $(1,\infty)$ is fixed by conjugation.  Its induced
isomorphism of geometric fibre functors identifies the two splittings.

The Hahn field has a henselian $t$-valuation with divisible value group
$\mathbb Q$ and residue field $\mathbb Q_p$.  The characterization in
Koenigsmann's \cite[\S1, Fact on $p$-adically closed fields]{B_Koenigsmann}
therefore makes $k$ $p$-adically closed.  In residue characteristic zero
there is no defect or wild inertia; divisibility of the value group
eliminates tame ramification.  Restriction and residue specialization
give $G_k\simeq G_{\mathbb Q_p}$.
Base change of the fundamental-group sequence of the constant curve
$U_0$ is consequently an isomorphism of extensions.  Henselian
specialization of finite \emph{\'etale} covers near $2$ identifies their
fibres at $2$ and $2+t$, compatibly with morphisms of covers and with
$G_k$.  This gives an isomorphism of fibre functors, and hence equality
of the section classes.  More explicitly, a simple root in a fibre at
$2$ has a unique lift in the prescribed residue class at $2+t$;
uniqueness makes these identifications compatible over the entire
finite-\'etale category.

An embedding of $k$ into a finitely generated extension $F/\mathbb Q_p$
would send $t\ne1$ to an element with an $n$th root for every $n$.
Such an element in $F$ has zero divisor on a normal proper model over
the algebraic closure of $\mathbb Q_p$ in $F$, hence lies in that finite
extension of $\mathbb Q_p$.  In a finite $p$-adic field
$F_0^\times\simeq\mathbb Z\times\mathcal O_{F_0}^\times$ and
$\bigcap_n(F_0^\times)^n=1$.  This contradicts injectivity of the
embedding.  This also supplies a direct verification of the required
``not sub-$p$-adic'' assertion.
\end{proof}

\subsection{The local construction and its tangential interpretation}

\begin{statement}{B.3}{Proved}{Uniformizers, tangent vectors, and Deligne's fibre functor}
Let $y\in X(k)$ and let $R$ be its henselian local ring, with fraction
field $L$ and residue field $k$ of characteristic zero.  A uniformizer
$t$ defines a splitting of $G_L\to G_k$.  Replacing $t$ by $t(1+h)$,
$h\in\mathfrak m_R$, preserves its inertia-conjugacy class.  The
splitting classes form a torsor under
$H^1(k,\widehat{\mathbb Z}(1))=\widehat{k^\times}$.
The classes defined by nonzero tangent vectors form the
$k^\times$-Kummer orbit in this torsor and are precisely Deligne's
tangential fibre functors.
\end{statement}
\begin{proof}
Choose compatible roots $t^{1/n}$ and use the compositum of the
unramified extension with these roots.  Since the residue characteristic
is zero, every finite extension becomes a Kummer extension of this
form after an unramified extension.  Acting on the coefficients and
fixing every $t^{1/n}$ gives the splitting.  Hensel's lemma gives unique
compatible roots in $1+\mathfrak m_R$ of $1+h$, so the displayed change
of uniformizer gives an isomorphic splitting.

Fix the inertia convention
$\zeta(t^{1/n})=\zeta_n t^{1/n}$ and
$\kappa(a)(\sigma)=\sigma(a^{1/n})/a^{1/n}$.
The splitting which fixes $(at)^{1/n}$ differs from the splitting which
fixes $t^{1/n}$ by $-\kappa(a)$.  On the other hand, a tangent vector $v$
with $t(v)=a$ differs from the vector with $t(v)=1$ by $+\kappa(a)$.
This explains the inverse transformation between a basis of
$\mathfrak m/\mathfrak m^2$ and its dual tangent vector; it also fixes
the sign rather than identifying these two torsors silently.

For an abelian kernel, differences of splittings are precisely continuous
$1$-cocycles and inertia conjugation changes a cocycle by a coboundary.
The Kummer sequence and Hilbert 90 identify this $H^1$ with the stated
inverse limit.  A finite cover near $y$, after an unramified extension,
has a parameter $u$ with $t=u^e$ times a unit.  Its fibre over a tangent
vector is the set of roots of the leading coefficient.  The same roots
are obtained by the preceding compatible-root construction.  This is
an isomorphism on every object and morphism in the finite-\'etale
category and is compatible with products.  It identifies the resulting
fibre functor with the specialization to the punctured tangent line
in Deligne \cite[\S15.13--15.26]{B_Deligne}.
Stix's local packet map \cite[Proposition 8 and Theorem 14]{B_StixCuspidal}
is induced by the same inclusion $G_L\to\Pi_U$; thus it is the packet
in the 1983 letter.  Over finitely generated fields of characteristic
zero that map is injective on splitting classes.
\end{proof}

\begin{statement}{B.4}{Proved}{Ramification and the two differential maps}
At rational boundary points of a finite morphism of curves, choose
parameters with
$t=a u^d(1+h)$, $a\in k^\times$, $h\in(u)$.
In tangent coordinates the induced map of packets is
\[
 \widehat{k^\times}\longrightarrow\widehat{k^\times},
 \qquad c\longmapsto a c^d.
\]
It restricts to the homogeneous map of tangent vectors
$v\mapsto a v^d$.  In the pure local model $t=a u^d$, the trace and
the determinant norm of differentials satisfy respectively
\[
 \operatorname{Tr}_{k((u))/k((t))}(du/u)=dt/t,
 \qquad
 \operatorname{Nm}(du)=
   (-1)^{d-1}d^{-d}a^{-1}t^{1-d}(dt)^{\otimes d}.
\]
The latter equality is in the norm of the one-dimensional differential
space, identified with $(\Omega^1_{k((t))/k})^{\otimes d}$ by $dt$.
\end{statement}
\begin{proof}
The factor $1+h$ has unique compatible roots in its principal-unit group.
On inertia the map is multiplication by $d$; on the chosen constant
splittings its discrepancy is the Kummer cocycle of $a$.  This proves
the first formula.  Equivalently the map on conormal leading terms is
$t\mapsto a u^{\otimes d}$, whose dual description is the displayed
homogeneous tangent map.

In the pure model, $du/u=d^{-1}dt/t$, and tracing the coefficient
$1/d$ multiplies by the field degree $d$.  The minimal polynomial of
$u$ gives $\operatorname{Nm}(u)=(-1)^{d-1}t/a$.  Since
$du=(ad u^{d-1})^{-1}dt$, taking determinants yields the norm formula.
The identity
$d\log\operatorname{Nm}(u)=\operatorname{Tr}(d\log u)$ follows at once.
For a nontrivial residue-field extension these formulas are applied
after splitting the residue points and then descended by the ordinary
field norm and cohomological corestriction.  Thus neither the trace
map nor the norm line is being substituted for the tangent map.
\end{proof}

\subsection{Every fibre of abelianization on points and packets}

\begin{statement}{B.5}{Proved}{Abelianized fibres in positive genus}
Let $k/\mathbb Q$ be finitely generated and let $g(X)\geq1$.
Write $r=|S(\bar k)|$.  In the map from rational points and full rational
cuspidal packets to geometrically abelianized section classes:
\begin{enumerate}
\item every rational point has a singleton fibre;
\item images belonging to different rational cusps are disjoint, and
      are disjoint from all rational point images;
\item for $r=1$ the entire unique packet has one abelianized image;
\item for $r\geq2$ the map on every packet is injective.
\end{enumerate}
For $r=0$ the domain consists only of the rational points.
\end{statement}
\begin{proof}
The geometric abelianization is the full Tate module of the generalized
Jacobian $J_S$, with the compact Jacobian $J$ as quotient.  After a
finite extension choose a base point.  The image in $H^1(k,TJ)$ of an
interior point or a packet centred at $z$ is the Kummer class of
$[z-b]$.  The Kummer injection of B.1, and the injectivity of
$X\to J$, separate distinct centres.  The latter injectivity follows
directly: a principal divisor $z-z'$ for $z\ne z'$ would define a
degree-one map to $\mathbb P^1$, contradicting $g\geq1$.
An equality before finite extension remains an equality afterwards,
so the reduction has lost no possible collision.

The toric part of $J_S$ has cocharacter module
$\mathbb Z[S(\bar k)]/\mathbb Z(1,\ldots,1)$.
For $r=1$ this module is zero: inertia maps to zero in
$\Delta_U^{\rm ab}$, proving (3).
For a rational cusp $y$ and $r\geq2$ its coordinate defines the exact
sequence of torus Tate modules
\[
 0\to\widehat{\mathbb Z}(1)\to TT_S\to TT_{S\setminus\{y\}}\to0.
\]
The two relevant degree-zero groups, $H^0(k,TT_{S\setminus\{y\}})$ and
$H^0(k,TJ)$, vanish.  Indeed they are compatible systems of rational
torsion, and the rational torsion groups of a torus or abelian variety
over a finitely generated characteristic-zero field are finite.
The two cohomology exact sequences therefore make
$H^1(k,\widehat{\mathbb Z}(1))\to H^1(k,TJ_S)$ injective.
This is exactly the difference map inside one packet.
\end{proof}

\begin{statement}{B.6}{Proved}{Genus-zero abelian fibres, including the exceptional collisions}
Let $k/\mathbb Q$ be finitely generated, $g(X)=0$, and put
$r=|S(\bar k)|$.  An anisotropic conic has neither rational interior
points nor rational cusps, so the domain in question is empty.
Otherwise choose a rational interior point as $\infty$ and write
$S\subset\mathbb A^1(\bar k)$.  In the coordinate torsor for
$H^1(k,\Delta_U^{\rm ab})$, the point and packet maps are represented by
\[
 v(x)_s=x-s,\qquad
 v(y,a)_s=\begin{cases}a&s=y,\\y-s&s\ne y,\end{cases}
 \quad a\in\widehat{k^\times},\ y\in S(k).
\]
Two such vectors represent the same class precisely when their
quotient is the restriction of a single diagonal element of
$\widehat{k^\times}$.
For $r\geq4$ every nonempty fibre on this domain is a singleton.  For $r=3$ the
only nonsingleton fibres consist of one element in each of two distinct
rational cusp packets.  For three rational cusps $y,z,w$, their unique
pairwise intersection is
\[
 a=\frac{y-w}{z-w}(z-y),\qquad
 b=\frac{y-z}{(y-w)/(z-w)}.
\]
The three pairwise intersections are distinct.  For $r=2$ with split
boundary, every abelian class has one element in each cusp packet and
additionally one interior point exactly when the class is in
$\kappa(k^\times)$.  For $r\leq1$ the geometric group is trivial and
all elements of the domain have the same image.
\end{statement}
\begin{proof}
Let $A=\Gamma(S,\mathcal O_S)$ and
$T_S=(\operatorname{Res}_{A/k}\mathbb G_m)/\mathbb G_m$.
For a character $(n_s)$ with $\sum_s n_s=0$, the generalized Abel map
evaluates to $\prod_s(x-s)^{n_s}$.  Specialization at $y$ replaces the
vanishing factor $x-y$ by the tangent parameter and the other factors
by $y-s$.  This proves the vector formulas, including their
$G_k$-equivariance for nonsplit $S$.
The Tate-module sequence of
$1\to\mathbb G_m\to\operatorname{Res}_{A/k}\mathbb G_m\to T_S\to1$
and Shapiro's lemma show that equality of these vector classes is
exactly equality modulo a diagonal Kummer class.  There is no extra
kernel here; the preceding term $H^0(k,TT_S)$ is zero.

One may check the equalities after a finite splitting field $L/k$.
The map $\widehat{k^\times}\to\widehat{L^\times}$ is injective:
the inflation--restriction kernel is
$H^1(\operatorname{Gal}(L/k),\widehat{\mathbb Z}(1)^{G_L})$,
and the coefficient group is zero because $L$ has only finitely
many roots of unity.  One first replaces $L$ by its Galois closure.
Two distinct nonvanishing entries of a point--cusp comparison would
require
$(x-s)/(y-s)=(x-s')/(y-s')$, hence $x=y$.  Thus no such collision
exists for $r\geq3$.  Point--point injectivity also follows from B.1.
For two distinct cusps, entries outside them require
$D=(y-s)/(z-s)$ for the same $D$ and every remaining $s$.
Two remaining entries force $y=z$; one remaining entry gives precisely
the displayed $a,b$.  The two intersection parameters within any one
packet differ by $-1$, whose full Kummer class is nonzero by B.1.
There is therefore no triple intersection.  The cases $r=2$ and
$r\leq1$ follow by the same formula from $T_S=\mathbb G_m$ and
$T_S=1$, respectively.  If only some boundary points are rational,
the asserted pairwise intersections occur only for those rational
packets; with $r=3$, two rational cusps force the third to be rational.
The finite rational-coordinate verification is recorded in
\texttt{code/B/packet\_and\_extension\_checks.py}.
\end{proof}

\subsection{Local sections and their limits}

\begin{statement}{B.7}{Proved using the prompt's local theorem}{The complete neighbourhood deduction of (T1)}
For a finite extension $k/\mathbb Q_p$ and a hyperbolic $U/k$, the map
\[
 U(k)\ \sqcup\!\coprod_{y\in S(k)}\operatorname{Sect}(D_y/G_k)/I_y
 \longrightarrow \operatorname{Sect}(\Pi_U/G_k)/\Delta_U
\]
is bijective.  Each packet is a torsor under $\widehat{k^\times}$.
The external input used for surjectivity is precisely Theorem 1.1 of
the attached October 6, 2026 preprint \cite{B_OpenAIPadic}; its
unrefereed provenance is retained in the source audit.
\end{statement}
\begin{proof}
For a section $s$, its neighbourhoods are the finite \emph{\'etale}
covers corresponding to open subgroups containing $s(G_k)$.
They have a cofinal subsystem whose smooth completions have genus at
least two.  To see this even in initial genus zero or one, first choose
a geometric finite \emph{\'etale} cover whose completion has genus at
least two, then choose a characteristic open subgroup $N\subseteq
\Delta_U$ contained in its subgroup.  The subgroup $N s(G_k)$ gives
the required neighbourhood.  Genus zero with at least three punctures
admits a Fermat-type cover; a once-punctured genus-one curve admits
the Heisenberg cover with nontrivial boundary inertia.  Riemann--Hurwitz
then gives the required genus.  Further finite covers preserve genus
at least two.

The image of the lifted section on each such completion is the section
of a unique $k$-point by the cited theorem.  These points are compatible:
a map of neighbourhoods sends one to a point inducing the other
section, and proper-curve injectivity gives equality.  Their common
image is a point $x\in X(k)$.  When $x\in U(k)$ the compatible lifts
give a rational pro-point in the section neighbourhood tower, and
the stabilizer of that pro-point is the section image.  When
$x\in S(k)$, they give a compatible point in the boundary of the
normalized tower, whose stabilizer is a cuspidal decomposition group.
The section therefore factors through that decomposition group.
This is the exact limit argument of
\cite[\S2.1.2]{B_StixLocal}.  The disjointness and injection of packets
are \cite[Theorem 14]{B_StixCuspidal}; their local Kummer description
is B.3.  Point injection follows from the local Jacobian Kummer map.
\end{proof}

\begin{statement}{B.8}{Proved}{The $p$-adic tangential subpacket and all its point limits}
For $k/\mathbb Q_p$ finite, a uniformizer $\varpi$ gives
\[
 k^\times=\varpi^{\mathbb Z}\times\mathcal O_k^\times,
 \qquad \widehat{k^\times}=\varpi^{\widehat{\mathbb Z}}
                         \times\mathcal O_k^\times.
\]
The Kummer map is injective, has cokernel
$\widehat{\mathbb Z}/\mathbb Z$, and has dense image.
At a rational cusp, every packet element is a limit of point sections
of points tending to that cusp in the $p$-adic topology.  Every
convergent limit of such point sections belongs to that packet.
\end{statement}
\begin{proof}
The unit group is a finite group times a finite-rank $\mathbb Z_p$
module; its power completion is itself.  The displayed formulas and
the injection, cokernel, and density follow.
Write a desired packet parameter as $(\alpha,u)$ with
$\alpha\in\widehat{\mathbb Z}$ and $u\in\mathcal O_k^\times$.
Choose increasing positive integers $n_j\to+\infty$ with
$n_j\equiv\alpha\pmod{j!}$, and put $t_j=\varpi^{n_j}u$.
A local \emph{\'etale} coordinate at the smooth rational cusp identifies
all sufficiently small $t_j$ with points $x_j\in U(k)$.
For any fixed finite quotient of the fundamental group, normalize
its cover at the cusp.  The local maps are Kummer maps times units,
and Hensel's lemma makes the unit factors have the required roots
on a sufficiently small disc.  Thus the point section modulo this
quotient is determined by the corresponding finite Kummer parameter.
The congruences on $n_j$ make it eventually the specified packet
parameter.  This proves convergence in every finite quotient.
Conversely the same description puts every sufficiently near point
in the image of the packet at every fixed finite level.  The packet
is compact and its map to the space of sections is continuous;
the compatible limit belongs to its image.  This proves the stated
equality of accumulation sets rather than only density of tangents.
\end{proof}

\begin{statement}{B.9}{Proved}{Genus zero and genus one over a $p$-adic field}
For a smooth proper genus-zero curve over $k/\mathbb Q_p$ there is
exactly one arithmetic section class, including for an anisotropic
conic.  For an elliptic curve $E/k$ with its origin, the section set is
$H^1(k,TE)$ and there is an exact sequence
\[
 0\longrightarrow E(k)\longrightarrow H^1(k,TE)
 \longrightarrow T H^1(k,E)\longrightarrow0.
\]
Local Tate duality identifies the last group with
$\mathbb Z_p^{[k:\mathbb Q_p]}$ noncanonically.
The point-section image is closed; its closure is exactly $E(k)$.
For a genus-one torsor $C$ under $E$, the section set is empty exactly
when its extension class $\delta([C])\in H^2(k,TE)$ is nonzero.
This class vanishes exactly for the maximal divisible subgroup of
$H^1(k,E)$.  When it vanishes, the section set is a torsor under
$H^1(k,TE)$.
\end{statement}
\begin{proof}
A genus-zero curve has trivial geometric fundamental group.
An explicit pointless conic is the projective norm equation
$N_{k'/k}(x)=\varpi z^2$ for an unramified quadratic $k'/k$:
the valuation of a nonzero norm is even, whereas the right side has
odd valuation.  Thus the unique section need not be a point section.

For $E$, the geometric fundamental group is $TE$ and the origin
splits the extension.  Continuous cocycles classify all splittings.
The inverse Kummer sequences give the exact sequence above;
the finite cohomology groups at every torsion level eliminate the
possible inverse-limit defect.  Since $E(k)$ is compact, its own
power completion is $E(k)$.  Local Tate duality makes $H^1(k,E)$
the Pontryagin dual of the compact group $E^\vee(k)$; this has a
finite group and a free $\mathbb Z_p$-module of rank $[k:\mathbb Q_p]$.
Taking Tate modules proves the formula.  The continuous point map
from the compact group $E(k)$ to the Hausdorff inverse-limit
cohomology group has closed image.

For a torsor, the finite-level extension classes are the connecting
classes of $[C]$ for multiplication by $n$ on $E$, as is seen by
identifying the automorphism group of $[n]:E\to C$ with the pulled-back
extension; this is also the cocycle-level calculation in
\cite[Proposition 2.2]{B_HarariSzamuely}.  All vanish exactly
when $[C]$ is divisible by all $n$.  The local-duality decomposition
just given identifies these elements with the maximal divisible
subgroup and supplies compatible divisions.  Finite-level
Mittag--Leffler then gives $\delta([C])=0$ in continuous degree two.
An abelian extension splits exactly when this class vanishes;
choosing a splitting identifies its section set with $H^1(k,TE)$.
Nonzero classes in the divisible subgroup give pointed-free genus-one
torsors with sections.  None is being used as a hyperbolic example.
\end{proof}

\subsection{Fixed groups, partial splittings, and the six axioms}

For reference, the actual axioms in scans 412--421 have the following
typed content.  This transcription specifies the quantifiers before
the geometric interpretation used below.
\begin{enumerate}
\item Admissibility of an outer action survives restriction to an open
subgroup of $\Gamma$.  The converse appears as a question in the margin.
\item Filling a $\Gamma$-stable set of holes preserves admissibility,
whenever the resulting type is still hyperbolic.
\item An admissible effective action has $\pi^\Gamma=1$.
\item An admissible outer action has an admissible effective lift
after restriction to some open subgroup.
\item For a fixed hole $i$, filling it maps admissible point classes
bijectively onto the admissible point classes of the filled group
other than the canonical class $P_i$.  Clause (a) excludes $P_i$,
clause (b) is injection, and clause (c) supplies the remaining classes.
The literal last sentence of scan 418 does not repeat ``other than
$P_i$''; without that exception it contradicts the displayed bijection
and clause (a).  The display is the formulation tested here.
\item Given a partial admissible action over an open subgroup,
and an action over the full group on the filled object preserving
the finite set of removed point germs, there is a unique admissible
extension of the action on the punctured object.  Scan 421's
normalizer calculation identifies its kernel with $L^{\Gamma'}$.
\end{enumerate}
The provisional definition of scans 431--432 imposes: (a) open image
in the arithmetic quotient $\boldsymbol\Gamma_a$; (b) no normalized
remaining inertia for the effective actions obtained by filling holes;
and (c) infinitely many conjugacy classes of section germs normalizing
no inertia.  Its asserted conclusions in scans 432--433 are:
(a) trivial fixed group for an effective admissible lift, with a
separate stronger question for $\ker\chi$; (b) the distinguished
effective action after filling one hole normalizes no remaining
inertia; (c) an arbitrary effective admissible action has the same
property after filling; and (d) the deletion map is injective with
complement exactly its distinguished point.  Scans 434--435 withdraw
the last surjectivity assertion for parameter spaces because distinct
scheme sections can meet.  B.16 realizes that issue by the actual
diagonal and constant graphs and separately refutes unrestricted
group-section surjectivity.

\begin{statement}{B.10}{Proved}{Every full-section action has trivial fixed group}
Let $k$ be finitely generated over $\mathbb Q$, or finite over
$\mathbb Q_p$.  For every section $s$ of a hyperbolic curve and every
open subgroup $G'\subseteq G_k$,
\[
 C_{\Delta_U}(s(G'))=1.
\]
In particular this condition, appearing as axiom (3) and in scan 431,
holds for point sections and for all cuspidal sections.
\end{statement}
\begin{proof}
Replace $k$ by the finite extension belonging to $G'$.  Suppose
$z\ne1$ is fixed.  Choose a characteristic open subgroup
$N\subseteq\Delta_U$ such that $z\notin N$.  The subgroup
$H=N\overline{\langle z\rangle}$ is normalized by $s(G_k)$.
It gives a section neighbourhood with geometric group $H$.
The image of $z$ in $H^{\rm ab}$ is nonzero, since its image in the
cyclic quotient $H/N$ is nonzero, and is fixed by $G_k$.
But $H^{\rm ab}$ is the Tate module of the generalized Jacobian of
that neighbourhood, whose rational torsion is finite in both field
classes under consideration.  Its invariant Tate module is zero.
This contradiction proves the assertion.  The use of
$N\overline{\langle z\rangle}$ is essential: a fixed element need not
already survive in the abelianization of $\Delta_U$ itself.
\end{proof}

\begin{statement}{B.11}{Proved}{The exact classification of effective lifts}
For a centre-free profinite geometric group $\Delta$, a continuous
outer action $\rho:G\to\operatorname{Out}(\Delta)$ has the canonical
extension
\[
 E_\rho=\operatorname{Aut}(\Delta)
             \times_{\operatorname{Out}(\Delta)}G.
\]
Its effective lifts are its splittings.  After choosing one lift
$\alpha$, their classes are exactly $H^1_{\rm cts}(G,\Delta_\alpha)$.
An outer action arising from a hyperbolic affine curve with a rational
cusp has effective lifts in that cusp packet.  Every such lift
normalizes the corresponding inertia group, including the
non-tangential elements of the packet.
\end{statement}
\begin{proof}
The kernel of $\operatorname{Aut}(\Delta)\to\operatorname{Out}(\Delta)$
is identified with $\Delta$ by the injective inner-action map.
For a second lift there is a unique $c_g\in\Delta$ with
$\alpha'_g=\operatorname{Inn}(c_g)\alpha_g$.  Multiplication of
automorphisms gives
$c_{gh}=c_g\alpha_g(c_h)$; conjugating the lift changes $c$ by a
coboundary.  These operations are continuous and are inverse to one
another.  For a cusp, the image is contained in $D_y=N_{\Pi_U}(I_y)$
by construction, independent of its Kummer parameter.  Thus the
whole packet consists of the pathological effective lifts excluded
by the no-normalized-loop condition, although their outer action is
the same geometric action as that of every rational-point lift.
Moreover $G_k^\circ=\ker\chi$ centralizes $I_y$, whereas B.10 proves
that the fixed group for an open subgroup of $G_k$ is trivial.
The cyclotomic-kernel statement for point sections is the
centralizer theorem proved in Part F.
\end{proof}

\begin{statement}{B.12}{Refuted}{Vanishing fixed kernel proves uniqueness, not existence}
For a partial splitting over an open normal subgroup $H\triangleleft G$
of an extension $1\to L\to E\to G\to1$, the equality $L^H=1$
does not imply that the partial splitting extends.  The exact
additional condition is that $N_E(s(H))\to G$ be surjective.
When it is surjective and $L^H=1$, its inverse is the unique extension.
\end{statement}
\begin{proof}
The kernel of this normalizer map is exactly $L^H$.
Surjectivity therefore makes it a continuous isomorphism of profinite
groups.  Every extending section lies in the normalizer, proving
uniqueness as well.

For a finite counterexample take
$L=S_3$, $H=S_3\times S_3$, and
$G=H\rtimes C_2$, where $C_2$ interchanges the factors.
Set $E=L\times G$ and $s(a,b)=(a,(a,b))$ over $H$.
The fixed subgroup in $L$ is $Z(S_3)=1$.
An extension would give a homomorphism $G\to S_3$ restricting to the
first projection on $H$.  The kernel of that restriction is
$1\times S_3$, which is not stable under interchange, a contradiction.
In fact the normalizer image is exactly $H$.
The supplied program enumerates all $432$ elements of $E$, verifies
the orders and the centralizer, and computes this normalizer image.
This refutes the implication proposed in the prompt from the
one fixed-kernel condition alone. The actual text of scan 421
already says that this condition implies uniqueness and that
existence means surjectivity of the normalizer projection.
The calculation therefore agrees with that distinction in
the manuscript. It does not delete the other geometric
hypotheses of axiom (6).
\end{proof}

\begin{statement}{B.13}{Proved}{The geometric meanings of axioms (1)--(6)}
For actions arising from pointed curves over the fields in B.10,
with admissible points defined as the actual interior point sections,
axioms (1)--(4) of scans 412--415 hold.  The three clauses (a)--(c) of
axiom (5) are exactly exclusion of the deleted point, injectivity,
and lifting of every other point.  They hold in the geometric
subcategory.  Axiom (6) holds there by effective descent of the
curve and its finite set of points, whenever the filled curve remains
hyperbolic.  The fixed subgroup of its actual forgetting kernel is
trivial, including when the removed points are unordered.  For
intrinsic admissibility defined only by absence of a
normalized cusp, axiom (5) is the additional point-surjectivity
assertion, after the already proved point and packet injections.
\end{statement}
\begin{proof}
Restriction to an open Galois subgroup is finite base change, giving
(1).  Filling a Galois-stable set of punctures is the open immersion
of curves into their partial completion, giving (2).  B.10 gives
(3).  A nonempty finite-type curve has a closed point with finite
residue extension, giving (4).
For a rational hole $y$, the actual point sets satisfy the literal
equality $(U\cup\{y\})(k)\setminus\{y\}=U(k)$; applying the
injective point-section map proves the three clauses of (5).
On all sections the omitted packet at $y$ maps to the point $y$,
so the set equality has exactly the meaning claimed only after the
admissibility condition is specified.

A finite set of distinct geometric points stable under descent data
descends as a finite \emph{\'etale} subscheme; deleting it commutes
with descent.  This gives existence in the geometric form of (6).
Here is a proof of the fixed-kernel assertion that also retains the
finite permutation quotient for an unordered removed set.  Write
$U\subset V$ for the punctured and filled curves, both hyperbolic,
and let $\ell$ in the kernel of the outer forgetting map be fixed
by the open subgroup $G_{k'}$.  Choose an automorphism $\beta$ of
$\Delta_U$ representing $\ell$.  The equality of outer classes
$[\beta]\rho(g)[\beta]^{-1}=\rho(g)$ gives the actual arithmetic
automorphism
\[
 (a,g)\longmapsto(\beta a\beta^{-1},g)
 \quad\hbox{of}\quad
 \operatorname{Aut}(\Delta_U)
       \times_{\operatorname{Out}(\Delta_U)}G_{k'}.
\]
B.11 identifies this pullback with $\Pi_{U_{k'}}$.
The relative Isom theorem, in the precise form of
Proposition~\ref{A:relative-hom}, realizes its outer class by a
unique $k'$-automorphism $f$ of $U$.  The preserved set of removed
cusps lets $f$ extend to $V$.  Its geometric outer class on $V$ is
the identity.  After an inner conjugation its arithmetic action
is the identity on $\Delta_V$; then its value at every element of
$\Pi_V$, divided by that element, lies in $\Delta_V$, centralizes it, and
is therefore the identity.  The injective Isom map for the
hyperbolic curve $V$ forces the extended morphism to be the
identity.  Restricting to $U$ gives $\ell=1$.
Thus the fixed kernel vanishes without replacing the unordered
forgetting kernel by its ordered subgroup.  The hyperbolicity of
$V$ is used in this Isom step.
The normalizer surjectivity furnished by descent is the extra step
which B.12 shows cannot be omitted.

For intrinsically selected sections, replacing point sections in this
argument by all sections which normalize no inertia group is exactly
the missing surjectivity assertion.  The cyclotomic centralizer
theorem identifies that set with the sections of zero
$G_k^\circ$-centralizer; it does not supply a pro-point.  B.7 supplies
the pro-point over $p$-adic local fields by its neighbourhood argument.
No corresponding unrestricted number-field pro-point construction
has been obtained here.
\end{proof}

\subsection{A global descent step beyond LM3 E.17}

\begin{statement}{B.14}{Proved}{Finite support at finite places forces a point or a cusp}
Let $U/k$ be hyperbolic over a number field and let $s$ be a Selmer
section.  Let $Z\subset X$ be finite over $k$.  Suppose the
compactification point of $s_v$ belongs to $Z(k_v)$ for a set of
finite places of density one.  Then $s$ is induced by a point of
$Z(k)\cap U(k)$ or is cuspidal at a point of $Z(k)\cap S(k)$.
In particular, a section whose finite localizations are all
cuspidal at one fixed rational cusp is globally cuspidal at that cusp.
\end{statement}
\begin{proof}
For a neighbourhood $V\to U$ with completion $Y$, let $Z_V$ be the
finite inverse image of $Z$.  The lifted section is again Selmer:
a lifted local point or cusp in the universal neighbourhood tower
projects to a local point or cusp in every intermediate cover.
The local compactification points lie in $Z_V$ at the prescribed
density-one set.  Its proper section gives an adelic point surviving
finite descent (pull back each finite torsor and use the section to
choose its arithmetic twist).  Passing cofinally to $Y$ of positive
genus, Stoll's Theorem 8.2 \cite{B_Stoll} identifies that adelic point
with a rational point of $Z_V$.  Thus $Z_V(k)$ is nonempty for every
neighbourhood.  These are finite sets.  Every finite collection of
compatibility conditions is realized by a point in a common
dominating neighbourhood, so compactness gives a compatible system
in $\varprojlim_V Z_V(k)$.

Its stabilizer in $\Pi_U$ contains the section image.  Over an
interior point this stabilizer is a point decomposition group; at
the boundary it is the cusp decomposition group.  The projection
to $G_k$ identifies the section as the respective point or
cuspidal splitting.  This is the argument in
\cite[\S3, proof of Corollary 5]{B_StixLocal}, with its final
noncuspidality assumption removed and the resulting cusp alternative
retained.  Only density-one finite places are used; no real
connected component has been replaced by a chosen point in the proof.
\end{proof}

\begin{statement}{B.15}{Proved using the prompt's local theorem}{Affine modular cuspidality}
For $N$ with $g(X_1(N))\geq2$, every section of
$\pi_1(Y_1(N))\to G_{\mathbb Q}$ is cuspidal on the affine curve.
Its centre is the unique rational cusp determined on $X_1(N)$.
\end{statement}
\begin{proof}
The compact modular conclusion is
\cite[Corollary 8.5]{B_OpenAIPadic}, combined with Mazur's torsion
theorem: $X_1(N)(\mathbb Q)$ consists of cusps at these levels.
This is the precise conclusion checked in LM3 D.10 and E.17.
Let $c$ be the resulting cusp.  At every finite place, B.7 makes
the affine localization a point section or a cusp section.
Its image on the compact curve is $s_c$; proper point injection
forces its centre to be $c$.  Thus it is cuspidal at $c$ at every
finite place.  B.14 now proves global cuspidality.
This last limit argument is the additional step absent from the
earlier return; equality on compactification alone was not used to
infer it.
\end{proof}

\subsection{Families and the obstruction to lifting a tautological section}

\begin{statement}{B.16}{Refuted}{Unrestricted section surjectivity for a constant family}
For $U_0=\mathbb P^1\setminus\{0,1,\infty\}$ over a number field $k$
and a smooth geometrically connected base $S/k$ with a rational
point, not every section of
$\pi_1(U_0\times S)\to\pi_1(S)$ comes from a morphism $S\to U_0$.
For geometric sections, deletion of a section retains precisely the
sections disjoint from it at every geometric point of the base.
\end{statement}
\begin{proof}
The product extension is
$\Pi_{U_0}\times_{G_k}\Pi_S\to\Pi_S$.
For a tangential cusp section $s_0:G_k\to\Pi_{U_0}$ at $0$, the map
$\gamma\mapsto(s_0(p_S(\gamma)),\gamma)$ is a splitting.
Its restriction at any $b\in S(k)$ is the cuspidal section $s_0$.
A graph $f:S\to U_0$ would restrict to the point section of $f(b)$,
contradicting the point--packet disjointness theorem.  This gives
an actual group-theoretic section and a complete counterexample.

For two morphism sections $f,g$, the graph of $f$ lies in the complement
of the graph of $g$ exactly when $f(b)\ne g(b)$ for every geometric
point $b$.  For $S=U_0$, the diagonal and the constant section $2$
meet at $2$ and are distinct, proving the correction to scans
433--435.  For the corresponding group-theoretic test, pull both
sections back along each closed-point section of $S$ and use
point-section injection over its finite residue field extension.
This tests precisely equality at that point.  The assertion is
about the indicated scheme sections; no unsupported lifting of a
group splitting through deletion is inferred.
\end{proof}

\begin{statement}{B.17}{Proved}{A tautological complete section cannot lift through its deletion}
Let $g\geq2$, $n\geq1$, and $k$ have characteristic zero.  For the
universal curve over $\mathcal M_{g,n/k}$, none of the $n$ tautological
complete sections admits a lift to a section of the fundamental-group
extension after all $n$ markings are removed.  This remains true
after every finite \emph{\'etale} base change of degree $d$.
The obstruction can be detected by a finite $\mu_m$-extension,
for a prime $m\nmid d(2g-2)$, with $d=1$ on the original base.
\end{statement}
\begin{proof}
Fix the proposed tautological marking, indexed by $1$.  Take a smooth
proper genus-$g$ curve $C/k$, put
$K=k(C^{n-1})$, and let $p_2,\ldots,p_n\in C(K)$ be its independent
generic points.  For $n=1$ put $K=k$ and omit this list.  After a
finite constant extension we may use a rational base point of $C$
to identify its Picard components with its Jacobian $J$.
Let $T=C_K\setminus\{p_2,\ldots,p_n\}$; varying the first marking
gives a morphism $T\to\mathcal M_{g,n/k}$.  The complete pulled-back
family is $C_K\times T$, and the first section is its diagonal.
Any supposed punctured lift, followed by filling the other holes,
would lift the diagonal homomorphism through
\[
 \pi_1((C_K\times T)\setminus\operatorname{diag}(T))
 \longrightarrow\pi_1(C_K\times T).
\]

We prove $c_1(T_C|_T)\ne0$ modulo $m$.  A contrary relation in
$\operatorname{Pic}(T)/m$ would give, on $C_K$,
\[
 T_C=mL+\sum_{i=2}^n a_i[p_i]
 \tag{B.1}\label{eq:B-Picrelation}
\]
in additive Picard notation.  Varying the generic tuple gives a
rational map $C^{n-1}\dashrightarrow\operatorname{Pic}^{\deg L}(C)$.
After translation it is a rational map to $J$, and extends to a
morphism: rational maps from a smooth variety to an abelian variety
extend across codimension one by properness and across the remaining
indeterminacy by the absence of rational curves in an abelian variety.
Its induced homomorphism on the $i$th Jacobian factor is an integral
endomorphism $\phi_i\in\operatorname{End}(J)$.
Equation \eqref{eq:B-Picrelation} gives
$m\phi_i=-a_i\operatorname{id}_J$.  Restriction to $J[m]$ shows
$m\mid a_i$.  Degrees in \eqref{eq:B-Picrelation} then imply
$m\mid(2g-2)$, a contradiction.  For $n=1$ this is simply the degree
argument.  Kummer's exact sequence injects
$\operatorname{Pic}(T)/m$ into $H^2_{\rm et}(T,\mu_m)$, proving the
claimed nonzero Chern class.

Let $\mathcal G$ be the $\mu_m$-gerbe of $m$th roots of
$\mathcal O(\operatorname{diag})$ on $C_K\times T$.
Its class is the Kummer Chern class.  On the complement of the
diagonal the canonical nonvanishing divisor section trivializes
this line, hence gives an actual morphism to $\mathcal G$.
The base $C_K\times T$ is an \emph{\'etale} $K(\pi,1)$ for finite
coefficients: its geometric factors are characteristic-zero
hyperbolic curves, their discrete surface/free groups are good, and
the geometric and arithmetic Hochschild--Serre sequences preserve
this comparison.  Thus the gerbe gives the genuine continuous
extension
\[
 1\to\mu_m\to\pi_1(\mathcal G)\to
                \pi_1(C_K\times T)\to1
\]
with the cyclotomic kernel action.  The asserted lift of the diagonal
would split its pullback.  That pullback class is
$c_1(\mathcal O(\operatorname{diag})|_{\operatorname{diag}})
=c_1(T_C|_T)$, just proved nonzero.  This contradiction proves the
assertion on the original base.

For a finite \emph{\'etale} moduli cover of degree $d$, pull it back
to the test curve $T$.  The resulting $f:T'\to T$ is finite
\emph{\'etale} of degree $d$, possibly disconnected.  On finite
coefficient cohomology its trace satisfies
$\operatorname{Tr}_f f^*=d$.  Choose $m\nmid d(2g-2)$.
The nonzero class $c_1(T_C|_T)$ therefore remains nonzero on $T'$;
otherwise its trace would give $d\,c_1(T_C|_T)=0$.
A lifted tautological section on the moduli cover would split the
gerbe pullback on every component of $T'$, contradicting this
nonvanishing.  The argument exhibits the extension with finite
kernel and its pullback; it does not require goodness of a mapping
class group.
\end{proof}

\begin{statement}{B.18}{Proved using the specified universal-curve input}{No sections of the universal punctured family in genus at least three}
Over $k/\mathbb Q$ finitely generated, for $g\geq3$ and $n\geq1$,
the universal punctured curve over $\mathcal M_{g,n/k}$ has no
fundamental-group section.  The external input is precisely the
complete-family section classification in
\cite[Main theorem]{B_LuoWatanabe}; its status and the inspected
proof step are recorded in the source audit.  The same conclusion
holds on its geometrically connected full symplectic level-$m$
covers over fields containing $\mu_m$.
\end{statement}
\begin{proof}
Fill the $n$ marked holes.  The cited classification makes the image
of any putative section one of the $n$ tautological complete classes.
Choose a representative and conjugate so that it is that tautological
section.  B.17 proves that this section has no lift through the
deleted family, giving a contradiction.  On a full symplectic
level cover the cited paper's final Remark gives the same complete
classification, and the finite-cover form of B.17 applies.
In particular the assertion
of scan 432 that there are no admissible sections holds here with
the stronger conclusion that there are no sections at all.
Calling a tautological complete section ``cuspidal'' after deleting
it does not by itself supply a splitting of its normal-line
extension; B.17 computes the obstruction to such a splitting.
\end{proof}

\subsection{Moduli criteria, lattices, and the additional archive tasks}

\begin{statement}{B.19}{Proved using the prompt's local theorem}{Local interpretation of admissibility and the stack $\mathcal M_{1,1}$}
For a hyperbolic curve over a finite extension of $\mathbb Q_p$, the
sections which normalize no cusp inertia subgroup are exactly its
point sections.  The same local point-or-cusp description, with its
full automorphism groupoids, holds for a hyperbolic DM curve, including
$\mathcal M_{1,1}$ and the $S_3$-quotient of the tripod.
\end{statement}
\begin{proof}
B.7 gives the two alternatives.  Every cusp alternative normalizes
its inertia by B.11, and point sections are disjoint from the
cuspidal packets.  Stix's normalizer characterization identifies a
section normalizing inertia with the corresponding packet.
For a DM curve choose a geometric finite \emph{\'etale} curve cover,
and a characteristic open subgroup $N$ of its geometric fundamental
group contained in the subgroup of that cover.  Finite generation
makes such $N$ available.  For the given section $s$, the subgroup
$Ns(G_k)$ defines a geometrically connected curve neighbourhood
$V\to\mathcal U$ over $k$ itself: its geometric group is $N$,
and its geometric cover is a scheme.  Descent of the finite
\emph{\'etale} covering retains that property.  B.7 applies to the
lifted section on $V$.  Pushing its point or its cusp fibre functor
forward supplies the required stack point or cusp class.
Full faithfulness is tested on these curve covers over finite
extensions: point injection and trivial point centralizers give it
on interior points; uniqueness of the boundary pro-point gives it
on cusp fibre functors and separates the two sorts of centres.
These are exactly the curve-neighbourhood lemma and four steps in
the proof of
\cite[Theorems B and C]{B_Steklov}, where Theorem C is the
full-faithfulness input over sub-$p$-adic fields.
Automorphisms, including the generic $\mu_2$ of
$\mathcal M_{1,1}$, are retained in this descent.  Passing just to
the coarse $j$-line would discard them and is not the map used.
\end{proof}

\begin{statement}{B.20}{Proved}{The lattice stabilizer for the tripod}
For the tripod over a number field $k$ with a rational base point,
fix an embedding into $\mathbb C$ and the resulting discrete
topological lattice in the geometric fundamental group.  Its
stabilizer in $G_k$ is trivial when the chosen embedding of $k$ is
nonreal, and is generated by its complex conjugation when that
embedding is real.  The oriented lattice has trivial stabilizer
in both cases.  Thus the canonical oriented discretization claimed
in scans 489--490 already fails for the tripod.
\end{statement}
\begin{proof}
Riemann existence maps the based discrete fundamental group injectively
and densely into its profinite completion.  A Galois automorphism
preserving this subgroup restricts to a discrete automorphism.
Its action on the discrete abelianization is the scalar $\chi$;
an integral invertible scalar has to be $+1$ or $-1$.  Each discrete
peripheral generator is therefore profinitely conjugate to the
corresponding discrete generator or its inverse.  Conjugacy
separability of a free group makes this an ordinary discrete
conjugacy.  Thus it preserves the three peripheral cyclic subgroups.
In the free group on two generators,
the pure peripheral outer automorphism group consists of the
identity and the orientation-reversing reflection: this is the
Dehn--Nielsen--Baer description of the extended mapping class group
of the three-punctured sphere.  Consequently its outer action is
either the identity or the outer action of complex conjugation.
Faithfulness of the tripod Galois action identifies the Galois
element itself in either case.  Complex conjugation belongs to
$G_k$ exactly for the chosen real embedding and preserves the based
lattice then.  It reverses its orientation, so cannot stabilize the
oriented lattice.  Changing the embedding by a Galois element
transports the comparison map and hence the lattice by that element;
the orbit is therefore nontrivial, as asserted.
\end{proof}

\begin{statement}{B.21}{Proved}{What the $\overline{\mathbb Q}$-algebraic entry supplies}
The reading of cote 148, scan 43 as obtaining ``sections Gal$/\mathbb Q$''
has the following verified pointwise content.  A finite-type
$\overline{\mathbb Q}$-curve with an algebraic point has, after a
finite field of definition, a point-induced section; its intrinsic
datum without that field is a germ over open subgroups of
$G_{\mathbb Q}$, modulo geometric conjugacy.  This datum determines
the point under the semi-abelian hypothesis of cote 143, scans 4--5.
The stronger reading asserting a $\mathbb Q$-model for every such
curve is false.

For a hyperbolic curve $U/\overline{\mathbb Q}$, the
constants condition in cote 143, scan 2 gives the exact sequence
\[
 1\longrightarrow\operatorname{Aut}_{\overline{\mathbb Q}}(U)
 \longrightarrow\operatorname{Aut}_{\rm abs}(U)
 \xrightarrow{\;\theta\;}\Gamma_U\longrightarrow1,
 \qquad
 \theta(f)=(f^*|_{\overline{\mathbb Q}})^{-1}.
\]
Its image $\Gamma_U$ is open in $G_{\mathbb Q}$, and need not
equal $G_{\mathbb Q}$. A finite field of definition gives a
continuous partial splitting. Conversely a continuous splitting
over $G_K$ gives an effective $K$-form of the curve, with its
boundary; conjugating by
$\operatorname{Aut}_{\overline{\mathbb Q}}(U)$ gives precisely
isomorphic forms.
\end{statement}
\begin{proof}
The equations of the curve, the chosen point, and finitely many
descent isomorphisms involve finitely many algebraic numbers, hence
are defined over one number field $L$.  The structure map of the
point gives a section over $G_L$.  Enlarging $L$ restricts the
section, which is exactly the germ convention in cote 143, scan 5,
and in the groupoid calculation of LM scans 61--65.
B.1 proves its injectivity on points.
For the additional exact sequence, the elements of
$\Gamma(U,\mathcal O_U)$ algebraic over $\mathbb Q$ are exactly
$\overline{\mathbb Q}$: a polynomial over $\mathbb Q$ splits
over $\overline{\mathbb Q}$, and the coordinate ring of the
integral curve has no zero divisors. An absolute automorphism
therefore preserves this constants subring. The inverse in
$\theta$ compensates for the contravariance of pullback, so
it is a homomorphism, with the asserted kernel and image.
A model over a number field $L$ gives the semilinear
base-change automorphisms with scalar pullback $\sigma^{-1}$,
for $\sigma\in G_L$. These form a section of $\theta$ over
$G_L$, proving that its image is open.

The geometric automorphism group of a hyperbolic curve is
finite. Enlarge $L$ so that the model and all its automorphisms
are defined over $L$. The inverse image of $G_L$ is then
$\operatorname{Aut}_{\overline{\mathbb Q}}(U)\times G_L$;
this specifies the natural profinite topology on the
absolute automorphism group. Two choices give the same topology:
their models and identifications are defined over one further
finite extension, on whose Galois group the model actions agree.
A continuous splitting over
$G_K$ differs from a model splitting, on a common open
subgroup, by a continuous cocycle with values in this finite
automorphism group. The equality locus of the two section
homomorphisms is an open subgroup. Its normal core in $G_K$
contains an open normal $G_M$, with $M/K$ finite Galois,
on which the cocycle is trivial. Normality makes every element
of the given $G_K$-action descend to an $M$-semilinear
automorphism of $U_M$; modulo $G_M$ these are the actual
$\operatorname{Gal}(M/K)$ descent datum on $U_M/M$.
It is effective: extend the automorphisms to the smooth
proper completion, take the invariant function field under
the finite semilinear group, and form its smooth proper
$K$-curve. Finite Galois descent of vector spaces identifies
the base-changed function field with the original one, hence
identifies the smooth proper models. The stable finite
boundary set descends and can be removed. Conversely every
$K$-form gives this datum, and an isomorphism of forms is
exactly the indicated conjugation. This proves the
effectivity statement in this precise hyperbolic-curve
range without treating an arbitrary abstract outer action
as an already effective descent datum.

For the stronger reading, take a pointed elliptic curve with
$j$-invariant $\sqrt2$.  Its $j$-invariant is preserved by every
isomorphism and is not rational, so it cannot descend to
$\mathbb Q$.  Its field-of-definition subgroup is nevertheless
open.  The doubtful two-word reading on scan 43 is retained as a
reading uncertainty; no conclusion depends on resolving it in
favour of the stronger assertion.
\end{proof}

\begin{statement}{B.22}{Proved}{Kummer structure, crossings, and finite covers in cote 143}
For a curve, a chosen lift of a boundary point gives an actual
homomorphism $\widehat{\mathbb Z}(1)\to\Delta_U$; changing that lift
conjugates the map.  Arithmetic transport acts through the
cyclotomic character and the permutation of boundary points.
At a strict normal crossing with $r$ branches, the corresponding
local group is $\widehat{\mathbb Z}(1)^r$ with its branch labels;
the actual map to the global geometric fundamental group must
also be retained.  It need not be injective.
For a finite \emph{\'etale} cover defined by $H\subseteq\Pi_U$,
the boundary lifts and inertias are
\[
 H\backslash\Pi_U/D_y,
 \qquad H\cap gI_yg^{-1},
\]
with the local ramification indices giving multiplication on each
one-dimensional Kummer lattice.  Every section germ lifts after
restriction to a sufficiently small open subgroup.
\end{statement}
\begin{proof}
The curve statement is the compatible-root construction of B.3
with a chosen point of the normalized universal cover.
At a strict normal crossing the local parameters
$t_1,\ldots,t_r$ may be rooted independently; Abhyankar's lemma
in characteristic zero gives the product local inertia.  Its
global map records which local Kummer coverings extend globally.
For example, remove two coordinate lines from $\mathbb P^2$.
The complement is $\mathbb A^1\times\mathbb G_m$.  At their
intersection use $x=X/Z$, $y=Y/Z$; the global multiplicative
coordinate is $u=y/x$.  Thus the local-to-global inertia map is
\[
 \widehat{\mathbb Z}(1)^2\longrightarrow\widehat{\mathbb Z}(1),
 \qquad (a,b)\longmapsto b-a,
\]
with nonzero diagonal kernel.  This provides the explicit
replacement for the dimension-one reservation in scans 6--8.

The double-coset and intersection formulas are the stabilizer
formulas for the action on the normalized universal cover.
Intersecting a procyclic inertia with an open subgroup gives
$e\widehat{\mathbb Z}(1)$, where $e$ is exactly the ramification
index.  For a section $s:G_L\to\Pi_U$, the inverse image
$s^{-1}(H)$ is open, and restriction lands in $H$; this is the
claimed lift of germs.  These are the actual maps behind
``automatiquement'' on cote 143, scan 12.  At crossings one uses
the corresponding finite-index sublattice and its integer
monomial map, rather than a single ramification integer.
\end{proof}

\begin{statement}{B.23}{Proved using the specified universal-curve input}{The unpointed universal curve in genus at least three}
For $k/\mathbb Q$ finitely generated and $g\geq3$, the fundamental-group
extension of the universal complete curve over $\mathcal M_{g/k}$ has
no section.  The same argument applies to its full symplectic level
covers included in \cite[Remark on topologically ample families]{B_LuoWatanabe}.
\end{statement}
\begin{proof}
Write $\mathcal C_g\to\mathcal M_g$ for this universal curve;
$\mathcal C_g=\mathcal M_{g,1}$.  Pull a proposed section back along
$\mathcal C_g\to\mathcal M_g$.  The resulting complete curve over
$\mathcal M_{g,1}$ has its diagonal marking.  The $n=1$ case of the
specified classification makes the pulled-back section conjugate
to that diagonal section.  Restrict to the geometric fibre $C$
of $\mathcal M_{g,1}\to\mathcal M_g$ over a fixed geometric curve.
In the product extension
$\Delta_C\times\Delta_C\to\Delta_C$, the section pulled back from
$\mathcal M_g$ is the graph of the trivial homomorphism, because
the fibre group maps to $1$ in $\pi_1(\mathcal M_g)$.
The diagonal is the graph of $\operatorname{id}_{\Delta_C}$.
Conjugation cannot identify a trivial homomorphism with the
identity on the nontrivial group $\Delta_C$.  This is the required
contradiction.  Marking a point does not alter a full symplectic
level structure, so the identical pullback diagram proves the
level assertion.
\end{proof}

\begin{statement}{B.24}{Proved}{Explicit low-genus finite \emph{\'etale} multisections}
The proposed absence of all admissible section germs for universal
families has the following geometric counterexamples.
\begin{enumerate}
\item Over $\mathcal M_{0,4}=\mathbb P^1_\lambda\setminus\{0,1,\infty\}$,
the universal four-punctured line acquires the interior section
$x=t$ after the degree-two \emph{\'etale} cover $\lambda=t^2$.
\item Over $\mathcal M_{1,1}$, the nonzero $m$-torsion of the universal
elliptic curve, for every $m\geq2$, is a finite \emph{\'etale}
multisection of the curve with its origin deleted.
\item Over $\mathcal M_2$, the six Weierstrass points are a degree-six
finite \emph{\'etale} multisection of the universal complete curve.
\end{enumerate}
\end{statement}
\begin{proof}
In (1) the covering curve is
$\mathbb P^1_t\setminus\{0,1,-1,\infty\}$.
Ramification of $t\mapsto t^2$ occurs only at $0,\infty$, already
removed.  The section value $t$ is different from $0,1,\infty,t^2$
on this curve: the only additional equality $t=t^2$ occurs at
$t=0,1$.  This is an explicit section of the pulled-back
\emph{punctured} family.
In (2), multiplication by $m$ on a characteristic-zero elliptic
scheme is finite \emph{\'etale}.  Removing its zero section from the
kernel leaves a finite \emph{\'etale} cover of the moduli stack of
degree $m^2-1$; the universal nonzero torsion point is the desired
section after pullback.
Every smooth genus-two curve is hyperelliptic.  Its canonical
degree-two map has six distinct simple ramification points in
characteristic zero, by Riemann--Hurwitz.  The canonical involution
and this reduced ramification divisor commute with base change;
the relative divisor is therefore finite \emph{\'etale} of degree
six over the smooth genus-two moduli stack.  Its universal point
gives (3).  Each construction supplies an open subgroup and an
actual algebraic section, hence its fundamental-group splitting.
The finite polynomial equalities used in (1) are checked in the
program supplied with this section.
\end{proof}

\begin{statement}{B.25}{Proved}{The necessary moduli tests and the exact tripod test}
An actual marked curve over a number field satisfies the distinctness
tests of scans 487--488: in positive genus the sections of its marked
points remain pairwise distinct after all holes are filled; in genus
zero they remain distinct in every three-puncture comparison in which
they are interior points.  Its rational interior point sections
normalize no boundary inertia.  For ordered type $(0,4)$ the moduli
point map and the tripod point map are literally the same map under
\[
 \mathcal M_{0,4}(k)=k\setminus\{0,1\},\qquad
 \mathcal M_{0,4}\simeq\mathbb P^1\setminus\{0,1,\infty\},
\]
where the four markings are $(0,1,\infty,\lambda)$.
The Isom-form reconstructs at most one marked curve from an outer
action which is already realized by a marked curve; it supplies no
point of this moduli curve for an arbitrary section.
\end{statement}
\begin{proof}
The positive-genus assertion is B.1 applied to the compact Jacobian.
For genus zero, retain three distinct marked points, send them by
the unique projective transformation to $0,1,\infty$, and apply B.1
to the resulting tripod.  The non-normalization assertion follows
from the inertia normalizer description and the disjointness of
point sections and cusp packets.  These are the actual quotient
maps involved in the manuscript tests.
The unique projective transformation also proves the displayed
moduli identification, over $k$ and on finite \emph{\'etale} covers.
Taking arithmetic fundamental groups and section classes therefore
identifies the two point maps, rather than merely comparing their
cardinalities.
Finally, two geometric realizations of the same peripheral outer
action give an isomorphism of their arithmetic curve groups over
$G_k$.  The Isom-form, with the marked inertia classes retained,
gives a unique isomorphism of marked curves.  Its input is an
isomorphism between the two realized curve groups; an arbitrary
moduli section contains no second realized curve to which that
theorem can be applied.  Thus the unresolved surjectivity here is
precisely the one recorded in B.27.
\end{proof}

\begin{statement}{B.26}{Proved}{The range of the Hom-form for constant families}
Let $S$ and $U$ be hyperbolic curves over a number field $k$.
A section of the constant-family extension is a homomorphism
$h:\Pi_S\to\Pi_U$ over $G_k$.  Every such $h$ with open image is
induced by a unique dominant morphism $S\to U$.
The homomorphisms with $h(\Delta_S)=1$ are exactly
$s\circ p_S$ for field sections $s:G_k\to\Pi_U$; among them, the
morphism-induced classes are exactly the constant morphisms at
$k$-points of $U$.  The corresponding graphs are the actual family
sections in these two ranges.
\end{statement}
\begin{proof}
The product fundamental group is the fibre product over $G_k$, so
projecting a splitting to the first factor gives $h$, and its graph
recovers the splitting.  Mochizuki's open Hom theorem applies to
$h$ with exactly its open-image hypothesis and produces the unique
dominant morphism.  Its graph is a section of the product family;
no claim that the graph itself is dominant in the surface is used.
For $h(\Delta_S)=1$, the quotient property of
$\Pi_S\to G_k$ gives a unique factor $s$, and being over $G_k$
makes $s$ a splitting.  A nonconstant morphism of curves is
dominant and induces an open image on geometric fundamental groups.
It therefore cannot induce this trivial geometric homomorphism.
The remaining morphisms are precisely the constant rational-point
maps.  B.16 supplies cuspidal examples among the other field-section
classes.  This proves the stated classification without applying
the dominant Hom theorem to a nondominant graph.
\end{proof}

\begin{statement}{B.27}{Open in this incomplete return}{The unresolved part of Part B}
The following requested conclusions have not been proved or
refuted here: point-surjectivity for all zero-cyclotomic-centralizer
sections of the tripod over a number field; the same point-or-cusp
dichotomy for every nonisotrivial hyperbolic curve over a finitely
generated positive-transcendence-degree extension of $\mathbb Q_p$;
all intrinsic instances of axioms
(5)--(6) outside the geometric subcategory; section classification
over arbitrary open subgroups of $\pi_1(\mathcal M_{g,n})$ and the
remaining cases in genus at most two; the general
intrinsic realization
criterion for arbitrary admissible outer actions.
\end{statement}
\noindent
The exact unresolved step for the first problem is construction
of an interior rational pro-point in the full section neighbourhood
tower.  B.10 and the cyclotomic-kernel theorem remove automorphisms
and identify the cusp alternative, but neither constructs that
pro-point.  B.14 constructs it for finite support; the unrestricted
tripod section does not acquire finite support from the argument.
The local theorem and B.7 yield local pro-points, whose compatibility
with one global rational pro-point is the missing descent step.
Passing to Jacobian Kummer classes determines a point already
present, but does not put every nonabelian section in the image.
For positive-transcendence-degree $p$-adic fields, B.38--B.40
now prove unramification and identify every proper generic
section with a relative fundamental-group section on a smooth
model. B.41 completes the isotrivial point-or-cusp classification
in every transcendence degree. For a nonisotrivial family the
remaining step is realization of its point-theoretic relative
section by a rational generic point. The factorization and
its fibrewise points alone do not prove that realization.

For arbitrary moduli covers, the inspected Luo--Watanabe statement
and its Remark on topologically ample families cover full symplectic
levels and bases inducing the stated geometric fundamental-group
isomorphism.  They do not quantify over every open subgroup.
Replacing an arbitrary open subgroup by a full level subgroup
would require a containment assertion not supplied by that result.
B.17 continues to exclude lifts of any tautological class, but
does not classify additional complete classes on such a cover.
The lattice question initially stopped at the discrete mapping
class group, but B.31 below completes it using the precise
archimedean intersection theorem.  The other gaps are recorded
as unfinished tasks, not as completed alternatives to the proof
requirements.

\subsection{Two exact repairs to the deletion statements}

\begin{statement}{B.28}{Refuted in the unrestricted formulation}{The collision test for equation (25)}
The map of scan 409, equation (25), is injective for a smooth
marked curve over a field finitely generated over $\mathbb Q$.
Its assertion for an arbitrary section of the arithmetic moduli
extension, without the distinctness admissibility condition of
scan 487, is false.  Already for $\mathcal M_{0,5}$ two different
marked holes can induce the same section after they are both filled.
\end{statement}
\begin{proof}
For a smooth marked curve the filled holes are distinct rational
points of the partially completed hyperbolic curve.  The point
injection of B.1 gives the assertion in this form, with no section
surjectivity hypothesis.

For the unrestricted assertion put $k=\mathbb Q$,
$R=k[[t]]$ and $K=k((t))$.  The ordered tuple
\[
 (0,1,\infty,2,2+t)
\]
defines a morphism $a:\operatorname{Spec}K\to\mathcal M_{0,5}$.
Let $r:G_k\to G_K$ be the uniformizer splitting of B.3.
Then $a_*\circ r$ is a section of
$\pi_1(\mathcal M_{0,5/k})\to G_k$.
The universal marked curve, pulled back by $a$, is the indicated
five-punctured projective line.  Fill its last two holes.  The
resulting curve is the constant tripod $U_0/k$, and the two
canonical sections of its arithmetic group are the composites
obtained from the $K$-points $2$ and $2+t$ with $r$.

Both points extend to $R$-points of $U_0$ with closed point $2$.
The specialization equivalence for finite \emph{\'etale} covers of
the henselian local scheme $\operatorname{Spec}R$ identifies
$\pi_1(R)$ with $G_k$.  In the commutative diagrams induced by the
two $R$-points, the map $G_k\xrightarrow{r}G_K\to\pi_1(R)$ is the
identity, and the map from the closed point to $U_0$ is $2$ in both
cases.  Thus both composites are $s_2$ up to geometric conjugacy.
They are the images of two different elements of $I'$, proving
failure of unrestricted injectivity.  In moduli this section is
tangential at the divisor where the fourth and fifth points
collide, so it fails exactly the distinctness test of B.25.
Scans 408--410 explicitly retain an admissibility reservation;
the example does not remove that reservation and then attribute
the enlarged assertion to a smooth marked curve.
\end{proof}

\begin{statement}{B.29}{Proved using the prompt's local theorem}{All fibres over the geometric locus of a deletion map}
Let $V/k$ be a hyperbolic curve over a number field, let $y\in V(k)$,
and put $U=V\setminus\{y\}$.  For the natural map
\[
 q_*:\operatorname{Sect}(\Pi_U/G_k)/\Delta_U
       \longrightarrow\operatorname{Sect}(\Pi_V/G_k)/\Delta_V,
\]
the fibre over $s_x$, $x\in V(k)\setminus\{y\}$, is the singleton
$\{s_x\}$; the fibre over $s_y$ is the entire cusp packet at $y$.
For every individual cuspidal class at a boundary point $z$ of $V$,
its fibre is one class in the corresponding packet of $U$.
Thus $q_*$ induces a bijection from the actual interior point
classes of $U$ onto those of $V$ other than $s_y$ and induces an
isomorphism on each old cusp packet.
For every proper hyperbolic $X/k$ with
$X(\mathbb A_k)_\bullet^{\mathrm{f\text{-}cov}}=X(k)$, all its
hyperbolic open subcurves have only point and cuspidal sections.
\end{statement}
\begin{proof}
Take a class in one of the specified fibres.  B.7 classifies its
localizations at every finite place; the real classification of
Part C supplies the remaining Selmer condition.  The corresponding local
classification for $V$ and the injection and disjointness of its
points and packets force every local compactification centre to
be $x$, $y$, or $z$, respectively.  B.14 with that single rational
point as finite support proves the same global alternative.
At $x\in U(k)$ there is just the point section.  At $y$ every
class of its packet maps to $s_y$, and no other class occurs.
At an old cusp $z$, the open immersion is an isomorphism on the
henselian punctured local neighbourhood.  Its homomorphism of
decomposition extensions identifies the two inertia groups and
their $\widehat{k^\times}$-torsors, proving the singleton assertion
for each specified cusp class.

For the final assertion, the stipulated proper theorem and its
finite-descent deduction make the compactification image of every
section the section of a rational point of $X$.  Proper point
injection at each finite place and B.7 force all local centres to
be that one point.  B.14 applies again.  This proves the advertised
affine conclusion on the indicated finite-descent class of proper
curves.

These fibre statements also give the exact logical use of axiom
(5): its clauses (a), (b), (c) are exclusion of $s_y$, injectivity,
and surjectivity onto the complement, on the admissible locus.
On a filled curve whose section set already consists of points
and packets, the computed fibres give the affine point-or-packet
classification and all three clauses.  For an arbitrary filled
curve, the remaining fibres are precisely those over section
classes outside its geometric locus; none is silently included
in this computation.
\end{proof}

\begin{statement}{B.30}{Proved using the prompt's local theorem}{A configuration section above distinct point classes}
Let $U/k$ be hyperbolic over a number field and let
$\operatorname{Conf}_n(U)\subset U^n$ be its ordered configuration
space.  Fix pairwise distinct $x_1,\ldots,x_n\in U(k)$.
The fibre of
\[
 \operatorname{Sect}(\pi_1(\operatorname{Conf}_n(U))/G_k)
       /\pi_1(\operatorname{Conf}_n(U)_{\bar k})
 \longrightarrow
 \prod_{i=1}^n\operatorname{Sect}(\Pi_U/G_k)/\Delta_U
\]
over $([s_{x_1}],\ldots,[s_{x_n}])$ is the singleton given by the
actual configuration $(x_1,\ldots,x_n)$.
\end{statement}
\begin{proof}
Proceed by induction on $n$.  The assertion for $n=1$ is literal.
Forgetting the last point is a hyperbolic-curve fibration
$p:\operatorname{Conf}_n(U)\to\operatorname{Conf}_{n-1}(U)$.
Its full geometric fundamental-group sequence is short exact,
including injection of the fibre group; this is
\cite[Proposition 2.2(i), with the total formation]{B_MochizukiTamagawa}.
All its hypotheses hold in characteristic zero.
By induction the image of the given section under $p_*$ is the
point section of $b=(x_1,\ldots,x_{n-1})$.
Conjugate the representative to have exactly that image.

The actual fibre is $U_b=U\setminus\{x_1,\ldots,x_{n-1}\}$.
The morphism $U_b\to\operatorname{Conf}_n(U)$ induces a continuous
isomorphism, with the fibre product taken along $s_b$,
\[
 \pi_1(U_b)\ \xrightarrow{\sim}\
 \pi_1(\operatorname{Conf}_n(U))
       \times_{\pi_1(\operatorname{Conf}_{n-1}(U))}G_k.
\]
Indeed it is the identity on the quotient $G_k$, and is an
isomorphism on geometric kernels by the homotopy exact sequence;
these facts give both injectivity and surjectivity.
Thus the given section is a section of $U_b/k$ whose image in
$\Pi_U$ is $s_{x_n}$.  Applying B.29 successively to the deleted
points makes its class uniquely the point $x_n$, since this point
is different from every deleted point.  In the configuration space
it is exactly $(x_1,\ldots,x_n)$.
Different conjugations of the base representative do not create
additional classes: the above argument supplies a geometric
conjugacy to the same actual configuration in every case.
\end{proof}

\begin{statement}{B.31}{Proved}{The lattice stabilizer in every hyperbolic type}
Let $U/k$ be hyperbolic over a number field, $x\in U(k)$, and
$\iota:k\hookrightarrow\mathbb C$.  Every extension
$j:\bar k\hookrightarrow\mathbb C$ gives a dense discrete subgroup
\[
 \Lambda_j=\operatorname{im}\bigl(
 \pi_1^{\rm top}(U_\iota(\mathbb C),x_\iota)
       \longrightarrow\Delta_U\bigr)
\]
whose profinite completion is $\Delta_U$.
Use $x$ for the effective action $\rho_x:G_k\to\operatorname{Aut}(\Delta_U)$,
and let $\sigma\cdot j=j\circ\sigma^{-1}$.  Then
\[
 \Lambda_{\sigma\cdot j}=\rho_x(\sigma)(\Lambda_j),\qquad
 \operatorname{Stab}_{G_k}(\Lambda_j)=G_{v,j}.
\]
Here $G_{v,j}=1$ at a complex place and is the order-two group of
complex conjugation at a real place.  The stabilizer of the oriented
lattice is $1$.  The unoriented orbit is $G_k/G_{v,j}$ and the
oriented orbit is a free transitive $G_k$-set of these comparisons.
Thus the effective action does not make the discretization
independent of the embedding choice.
\end{statement}
\begin{proof}
Riemann existence and residual finiteness of surface and free
groups give the injection and completion assertion.  Applying
$j\circ\sigma^{-1}$ to a finite cover and its fibre at $x$ gives
the displayed equivariance, including its inverse convention.

The exact arithmetic input is
\cite[Corollary 2.25]{B_HoshiMochizukiIV}:
inside $\operatorname{Out}(\Delta_U)$, the intersection of the
Galois image with
$\operatorname{Out}(\pi_1^{\rm top}(U_\iota(\mathbb C)))$ is $G_{v,j}$.
An automorphism preserving the actual based subgroup $\Lambda_j$
restricts to its discrete automorphism, so its outer class belongs
to that intersection.  The stabilizer is therefore contained in
$G_{v,j}$.  At a real place, complex conjugation is an actual
homeomorphism fixing $x_\iota$, and hence preserves the based lattice.
This proves the reverse inclusion.  It reverses the surface
orientation, giving the oriented assertion.  Orbit--stabilizer
and equivariance give the two orbit descriptions.

The cited theorem's proof was checked at the needed bridge:
its cartesian discrete/profinite configuration square
(Corollary 2.20(v)) lifts the discrete outer action to the third
configuration space; compatible tripod homomorphisms then reduce
the intersection to the faithful arithmetic tripod action.
This supplies the general-type step absent from the elementary
tripod proof B.20.
\end{proof}

\begin{statement}{B.32}{Proved using the prompt's local theorem}{All sections of proper constant families over a finite $p$-adic field}
Let $k$ be finite over $\mathbb Q_p$, let $S/k$ be a
hyperbolic curve, and let $X/k$ be a smooth proper
geometrically connected curve of genus at least two.
Every continuous homomorphism
$h:\Pi_S\to\Pi_X$ over $G_k$, modulo conjugation by
$\Delta_X$, is induced by a unique $k$-morphism $S\to X$.
Consequently the graph construction is a bijection
\[
 \operatorname{Hom}_k(S,X)
 \ \xrightarrow{\ \sim\ }\
 \operatorname{Sec}\bigl(\pi_1(X\times_k S)\to\Pi_S\bigr),
\tag{B.32.1}
\]
where section classes on the right are taken modulo
the geometric fibre group $\Delta_X$.
This closes the intermediate-image case of B.26 for
this field and proper-fibre range.
\end{statement}
\begin{proof}
For each finite extension $L/k$ and $a\in S(L)$,
the composite of the point section of $a$ with $h$
is a section of $\Pi_{X_L}\to G_L$.
The prompt's local theorem, used in B.7 with its
stated provenance, gives a unique point
$b\in X(L)$ inducing this composite. Thus $h$ is
point-theoretic in the exact sense of
\cite[Definition 3.2]{B_CarlsonStix}: it preserves
geometric section germs over all finite extensions.
Here properness excludes any cusp alternative.

The field $k$ is Kummer-faithful. For completeness,
an abelian variety's group $A(k)$ has an open
subgroup isomorphic to $\mathbb Z_p^d$ and finite
quotient, hence has no nonzero element divisible
by every positive integer. A split torus has the
same property from
$k^\times\simeq\mathbb Z\times\mathcal O_k^\times$;
a finite splitting extension treats all tori.
For an extension $1\to T\to G\to A\to1$, the
argument of B.1 applies with the finite exponent
of $A(k)_{\mathrm{tors}}$, and then that of
$G(k)_{\mathrm{tors}}$, to give
$\bigcap_n nG(k)=0$. This is the required
Kummer-faithfulness, and the same proof applies
over every finite extension.

Carlson--Stix, Proposition 3.17
\cite{B_CarlsonStix}, now gives the precise
dichotomy: $h$ is open, or its image is contained
in one point decomposition group. We use the
hyperbolic-target range of that proposition.
Its proof was checked at the use of arbitrarily
large finite-\'etale target covers, the
Abel--Jacobi dimension estimate of Proposition 3.16,
and the normalizer argument.
In the open case, the open Hom theorem of
Mochizuki, as in Proposition~\ref{A:relative-hom}, gives a unique dominant
$k$-morphism $S\to X$ inducing $h$.

In the other case the induced map on geometric
points is constant. Its equivariance under
$G_k$ makes the image a point $x\in X(k)$.
The image of $h$ is contained in a decomposition
group $D_x$; since both $h$ and $D_x\to G_k$
are over $G_k$, and the latter map is an
isomorphism, necessarily
$h=s_x\circ p_S$. This is induced by the
constant morphism at $x$.
Uniqueness of constant morphisms follows from
point-section injectivity. A dominant morphism
has open image, whereas $D_x$ has infinite
index, so the two cases do not overlap.

Finally, the product homomorphism is the actual
isomorphism of arithmetic extensions
\[
 \pi_1(X\times_k S)
       \simeq \Pi_X\times_{G_k}\Pi_S.
\]
On geometric kernels it is the product
isomorphism from the finite-\'etale K\"unneth
theorem, and on arithmetic quotients it is the
identity. A splitting of the projection to
$\Pi_S$ is therefore exactly
$\gamma\mapsto(h(\gamma),\gamma)$.
For a morphism $f:S\to X$ this is the induced
homomorphism of its graph
$a\mapsto(f(a),a)$, and geometric-fibre
conjugation is exactly the indicated conjugation
of $h$. These inverse constructions prove
(B.32.1).
\end{proof}

\begin{statement}{B.33}{Proved using the prompt's local theorem}{The complete constant-family decomposition and its cusp parameters}
Let $k$ be a number field or a finite extension of
$\mathbb Q_p$, and let $S/k$ and $U=X\setminus D$ be
hyperbolic curves. Write
$\operatorname{Sec}^{\mathrm{nc}}(U/k)$ for the field
section classes outside all cusp packets.
For each $y\in D(k)$ choose its inertia and decomposition
groups $I_y\subset D_y$ and put
\[
 \mathcal C_y(S)=
 \operatorname{Hom}_{G_k}(\Pi_S,D_y)/I_y.
\]
There is a disjoint decomposition, through the actual
product and graph maps of (B.32.1),
\[
 \operatorname{Sec}\bigl(\pi_1(U\times_k S)\to\Pi_S\bigr)
 \simeq
 \operatorname{Hom}^{\mathrm{dom}}_k(S,U)
 \ \amalg\
 \operatorname{Sec}^{\mathrm{nc}}(U/k)
 \ \amalg\
 \coprod_{y\in D(k)}\mathcal C_y(S).
\tag{B.33.1}
\]
Each $\mathcal C_y(S)$ is a torsor under
$H^1_{\mathrm{cts}}(\Pi_S,\widehat{\mathbb Z}(1))$.
The choice of a nonzero tangent at $y$ supplies its
origin. If $k$ is a number field, there is the canonical
Kummer identification
\[
 H^1_{\mathrm{cts}}(\Pi_S,\widehat{\mathbb Z}(1))
 \simeq
 \varprojlim_n
  \mathcal O(S)^\times/
                      \mathcal O(S)^{\times n}.
\tag{B.33.2}
\]
The right side is a power completion and is not
asserted compact. Inflation from $G_k$ is exactly
the subgroup arising from the constant units
$k^\times$. Thus all intermediate geometric images
in a constant family are accounted for by these
cusp torsors. The remaining noncuspidal
surjectivity question in (B.33.1) is precisely the
field-section question already isolated in B.27.
\end{statement}
\begin{proof}
The graph isomorphism identifies the left side with
$\operatorname{Hom}_{G_k}(\Pi_S,\Pi_U)/\Delta_U$.
Apply the infinite-index trichotomy of
Proposition~\ref{F-local-dichotomy}, including its
finite-$p$-adic proof. Its geometric image is either
open, zero, or a nonzero subgroup of one $I_y$.
In the first case Proposition~\ref{A:relative-hom}
gives exactly the dominant morphisms in (B.33.1).
In the second case the quotient property of
$\Pi_S\to G_k$ gives a unique field-section
factorization. Its noncuspidal classes give the
middle summand; its cuspidal classes belong to the
last summands. In the third case the inertia
normalizer identity puts the entire image in
$D_y$, and projection onto $G_k$ makes $y$
$k$-rational. This proves exhaustion.

Choose a tangential splitting $s_y:G_k\to D_y$.
Every homomorphism into $D_y$ has a unique expression
\[
 h(\gamma)=c(\gamma)s_y(p_S(\gamma)).
\]
The homomorphism equation is the continuous
$1$-cocycle equation for $c:\Pi_S\to I_y$, where
the action is cyclotomic through $G_k$.
Conjugation by $I_y$ changes $c$ by a coboundary,
giving the stated torsor.
We check that passage from $I_y$-conjugacy to
$\Delta_U$-conjugacy makes no additional
identifications. For a nonzero geometric image,
the peripheral intersection theorem forces a
geometric conjugator between two maps with image
in $D_y$ to belong to $I_y$. This is the same
nonzero-subgroup inertia calculation used in
Proposition~\ref{F-local-dichotomy}.
For zero geometric image the assertion is exactly
injectivity of the field cusp packet. Different
cusps cannot meet in either case, and zero and
nonzero geometric image are preserved under
conjugation. Finally the open case has infinite
nonprocyclic geometric image, so it meets none
of these summands. The field-section summand
was defined by removing its cusp packets.
This proves disjointness and injectivity.

For (B.33.2), the Kummer sequence on $S$ gives
\[
 0\to\mathcal O(S)^\times/\mathcal O(S)^{\times n}
   \to H^1(S,\mu_n)\to\operatorname{Pic}(S)[n]\to0.
\]
The finite-coefficient group in the middle is
$H^1(\Pi_S,\mu_n)$ by the classification of
finite-\'etale torsors. The degree map on the
smooth proper completion of $S$ has image in
$\mathbb Z$ and kernel a subgroup of its
Jacobian's $k$-points. Mordell--Weil therefore
makes that completion's Picard group finitely
generated. Quotienting by the finitely many
boundary-divisor classes gives
$\operatorname{Pic}(S)$, also finitely generated.
Consequently $T\operatorname{Pic}(S)=0$.
The inverse system of unit quotients is
surjective, while the finite groups
$H^0(S,\mu_n)$ are Mittag--Leffler along a
factorial tower. Taking inverse limits in the
displayed sequence and in continuous degree-one
cohomology therefore proves (B.33.2).

Naturality of Kummer identifies inflation with
constant units. More explicitly,
$\mathcal O(S)^\times/k^\times$ embeds by its
boundary divisor into a finitely generated
free abelian group, so it is free of finite
rank. The resulting unit sequence splits as
an abstract abelian-group sequence; its power
completions remain exact and the constant
subgroup injects. The kernel of restriction
from $H^1(\Pi_S,I_y)$ to $H^1(\Delta_S,I_y)$
is exactly the inflated $H^1(G_k,I_y)$ by
the continuous Hochschild--Serre sequence.
Hence the maps in this subgroup are precisely
the field cusp sections composed with $p_S$.
All other classes in the relative cusp torsor
have the nonzero intermediate image already
described. A morphism-induced constant class
in the middle summand is exactly an actual
point of $U(k)$, by B.26. This proves the final
assertion without filling the remaining field
section term by an unproved point set.
\end{proof}

\begin{statement}{B.34}{Refuted}{The unrestricted interior-point reflection in the printed RNS Theorem F(ii)}
Let $k=\mathbb Q_5$, and let
\[
 Y=\mathbb P^1_k\setminus\{0,1,\infty,2\},\qquad
 Z=\mathbb P^1_k\setminus\{0,1,\infty\},\qquad
 f:Y\hookrightarrow Z.
\]
A tangential section $s$ at the cusp $2$ of $Y$ is not a
point section of $Y$, whereas $f_*s$ is the point section of
$2\in Z(k)$. This satisfies every field, prime, subgroup and
map hypothesis in the literal interior-point statement of
Mochizuki--Tsujimura, Theorem F(ii), and Proposition 3.5(iii)
\cite{A_MochizukiTsujimuraRNS}. Thus the stated equivalence for
arbitrary dominant morphisms, with point sets $Y(\Omega^H)$
and $Z(\Omega^H)$, is false. The same curves and sections
over $\mathbb Q$ refute its unrestricted number-field analogue.
This statement concerns that specific reflection assertion;
it makes no assertion against the source's RNS Theorem A or
its Isom theorems.
\end{statement}
\begin{proof}
Choose $t=x-2$ and a nonzero tangent at $2$. The compatible-root
construction of B.3 gives
\[
 s:G_k\longrightarrow D_2(Y)\longrightarrow\Pi_Y,
 \qquad
 1\longrightarrow\widehat{\mathbb Z}(1)
 \longrightarrow D_2(Y)\longrightarrow G_k\longrightarrow1.
\]
The inclusion $f$ extends across $t=0$. It kills this inertia,
and its restriction to $D_2(Y)$ factors through the residue
map $D_2(Y)\to G_k$ followed by the point decomposition
section at $2$ in $Z$. Hence $f_*s=s_2$ up to geometric
conjugacy, for every tangent and indeed every member of
the packet.

The invertible function $x$ on $Y$ gives the actual map
$x_*:\Pi_Y\to\Pi_{\mathbb G_m}$.
At $2$, write $x=2(1+t/2)$; the factor $1+t/2$ has unique
compatible roots with residue $1$. Therefore $x_*s$ has
Kummer class $\kappa_k(2)$. If $s$ were induced by
$y\in Y(k)$, functoriality would give
$\kappa_k(x(y))=\kappa_k(2)$.
The local Kummer injection proved in B.8 forces $x(y)=2$,
which is excluded from $Y$. Over $\mathbb Q$, the same
argument uses B.1, or directly the valuations and the
finite group $\{\pm1\}$ to show
$\bigcap_n(\mathbb Q^\times)^n=\{1\}$.

Here are the exact local hypotheses. Take $K=\mathbb Q_5$,
$p=5$, $\Sigma$ equal to all primes, $\ell=3$, and
$H=G_K$. The residue field is $\mathbb F_5$; both curves
and the dominant open immersion are defined over the
$p$-adic local field $K$ itself. Tame inertia gives a
continuous surjection $H\cap I_K=I_K\twoheadrightarrow
\mathbb Z_3(1)$, hence to $\mathbb Z_3$ as a profinite
group. The $3$-adic cyclotomic image is the closure of
$5^{\mathbb Z}$ in $\mathbb Z_3^\times$, which is open:
$v_3(5^2-1)=1$, so the powers of $25$ generate
$1+3\mathbb Z_3$. Finally
$\Omega=\mathbb C_5$ and $\Omega^H=\mathbb Q_5$.
For the last equality one may apply the Ax--Sen--Tate
estimate recalled and sharpened in \cite{B_LeBorgne}:
an invariant element can be approximated arbitrarily
closely by elements of the complete field $K$ and
therefore belongs to $K$.

The source's printed page 8 uses the actual point sets
$Y(\Omega^H)$ and $Z(\Omega^H)$; its scheme and curve
conventions on pages 13--14 do not replace them with
compactifications. No finiteness or monocuspidal hypothesis
is imposed on $f$ there. In the proof of Proposition
3.5(iii), printed page 79, the reduction replaces finite
section neighbourhoods by their smooth compactifications.
In this example, every such compactification contains
the centre of the lifted cusp. The resulting point
does not lie in the original open neighbourhood.
That reduction therefore does not preserve the
interior-point property needed for the claimed equivalence.
\end{proof}

\begin{statement}{B.35}{Proved, using the stipulated local theorem}{The exact fibre over a point under a dominant morphism}
Let $k$ be a number field or a finite extension of
$\mathbb Q_p$, and let $f:Y\to Z$ be a dominant morphism
of hyperbolic curves over $k$. Write
$\bar f:\bar Y\to\bar Z$ for its extension to smooth
proper completions. For $z\in Z(k)$, the full-section
fibre is the disjoint union
\[
 f_*^{-1}([s_z])
 =\{[s_y]:y\in Y(k),\ f(y)=z\}
 \ \sqcup\!
 \coprod_{\substack{c\in(\bar Y\setminus Y)(k)\\
                         \bar f(c)=z}}
                  \operatorname{Packet}_c(Y).
 \tag{B.35.1}
\]
In particular, dominant morphisms preserve point
geometricity, and they reflect it on noncuspidal source
sections. Finite morphisms preserve and reflect point
geometricity on all sections. For finite \emph{\'etale}
morphisms the latter assertion follows directly from
the finite-cover action on the fibre over $z$, without
the stipulated local theorem and over every
characteristic-zero field.
\end{statement}
\begin{proof}
The map $\bar f$ is finite. A point section at $y$ maps
to that at $f(y)$. For a boundary point $c$ with
$\bar f(c)=z\in Z(k)$, the henselian map extends into
$Z$ at $c$, so its cusp inertia is killed and the
entire packet maps to $s_z$, exactly as in B.34.
These are all the displayed inclusions.

Over a finite extension of $\mathbb Q_p$, B.7 classifies
every source section as a point or a cusp. If a cusp
maps to a target cusp, B.4 maps its packet into that
target packet, which is disjoint from the point
sections. The remaining centres map into $Z$; local
point-section injection then makes their image equal
to $z$. This proves (B.35.1) locally.

Now let $k$ be a number field and let $s$ belong to
the left side of (B.35.1). At every finite place,
the local result just proved puts the compactification
centre of $s_v$ in the fixed finite $k$-subscheme
\[
 W=(\bar f^{-1}(z))_{\mathrm{red}}\subset\bar Y.
\]
The stipulated local theorem, B.7, supplies all finite
Selmer localizations; the real section comparison
of C.4 supplies the real ones. Apply B.14 to this
single finite support $W$. It proves that $s$ is a
point section at a member of $W(k)\cap Y(k)$ or a
cusp section at a member of
$W(k)\cap(\bar Y\setminus Y)(k)$.
Full point/packet injection and disjointness give
exactly the disjoint union in (B.35.1).

If $f$ is finite, normalization in $k(Y)$ identifies
$Y$ with $\bar f^{-1}(Z)$. Thus no boundary point
$c$ lies over $z\in Z$, and the packet terms vanish.
For finite \emph{\'etale} $f$, there is also the direct
proof valid over an arbitrary characteristic-zero
field. Identify $\Pi_Y$ with the corresponding open
subgroup of $\Pi_Z$. After geometric conjugation,
a section mapping to $s_z$ makes the appropriate
coset in the geometric finite fibre over $z$ fixed
by $s_z(G_k)$. That fixed coset is an actual
$k$-point $y$ of the finite \emph{\'etale} scheme
$f^{-1}(z)$. The decomposition group at that point
identifies the source section with $s_y$ inside
$\Pi_Y$. Conversely $s_y$ maps to $s_z$ by
functoriality. This proves the final assertion
through the actual finite-cover map.
\end{proof}

\subsection{The closed obstruction supplied by stable boundary cycles}

\begin{statement}{B.36}{Proved}{A closed obstruction to smooth arithmetic realization}
Fix a hyperbolic type $(g,n)$, its geometric full profinite
curve group $\Delta$, and its prescribed peripheral conjugacy
classes. Let $Q\subset\operatorname{Out}(\Delta)$ be the closure
of the geometric mapping class group, and let
$A\subset\operatorname{Aut}(\Delta)$ be its inverse image.
For each nonempty stable multicurve type, choose oriented
generators $c_1,\ldots,c_r\in\Delta$ of its node groups. Call
their $A$-translates, with arbitrary individual conjugations,
\emph{pro-simple multicurve tuples}.

Let $k$ be a number field, let $E\twoheadrightarrow G_k$ be a
profinite extension, and let
$\rho:E\to\operatorname{Out}(\Delta)$ be continuous, with the
image of the geometric kernel contained in $Q$. The following
condition on a continuous section $s$ defines a closed,
geometric-conjugacy-invariant subset of its section space:
there is a nonempty pro-simple multicurve tuple such that,
for every $\gamma\in G_k$, there are a permutation $\sigma$
and signs $\epsilon_i\in\{1,-1\}$ satisfying
\[
 \rho(s(\gamma))(c_i)
 \ \sim_\Delta\quad
 c_{\sigma(i)}^{\epsilon_i\chi(\gamma)}
 \qquad(1\leq i\leq r).
 \tag{B.36.1}
\]
Here the left side is a conjugacy class, so the condition is
independent of the representative of the outer automorphism.
No outer action satisfying this condition is the peripheral
outer action of a smooth marked curve over $k$, or over any
finite extension of $k$ after restriction.
\end{statement}
\begin{proof}
We first make the two compactness and nonperipherality
assertions precise. A finitely generated profinite group has
a basis of open characteristic subgroups; its automorphism
group is the closed inverse limit of the compatible finite
quotient automorphism groups. Hence $\operatorname{Aut}(\Delta)$
and its closed subgroup $A$ are compact. For a fixed tuple the
image of
\[
 A\times\Delta^r\longrightarrow\Delta^r,\qquad
 (a,h_1,\ldots,h_r)\longmapsto
 (h_i a(c_i)h_i^{-1})_{i=1}^r
 \tag{B.36.2}
\]
is compact. There are only finitely many stable multicurve
types: the number of curves is at most $3g-3+n$, and cutting
along them records a finite graph, the genera of its vertices,
and the distribution of the $n$ labelled punctures. Stability
and the genus formula bound all of these data. Classification
of the cut oriented surfaces identifies the resulting finite
list of mapping class orbits. Thus the disjoint union
$\mathcal K$ of the compact tuple spaces (B.36.2), over that
finite list, is compact. For $(g,n)=(0,3)$ the list is empty
and the assertion is vacuous.

Each initial $\langle c_i\rangle$ is an injected copy of
$\widehat{\mathbb Z}$, and is not contained in a peripheral
group. One elementary finite-cover verification is useful
here. A nonseparating curve has a primitive nonzero class
in the homology of the filled surface. For a separating
essential curve, each cut side admits a connected double
cover trivial on that boundary circle: on a positive-genus
side use a nonzero handle character to $\mathbb Z/2$; on
a genus-zero side use the character nonzero on two of its
punctures. The latter side has at least two punctures because
the curve is essential and nonperipheral. Gluing the two
connected covers along their two boundary lifts gives a
double cover in which each lift of the original curve is
nonseparating, including after the punctures are filled.
Its class is therefore primitive in compact homology.
Riemann existence gives the same finite cover of the
profinite surface group. A cusp subgroup, and its intersection
with this finite-cover subgroup, maps to zero in compact
homology. This proves the claimed injection and exclusion.
The groups belonging to distinct entries of a multicurve
are pairwise nonconjugate, by the distinct-edge assertion
in \cite[Proposition 1.2]{F-MochizukiCombinatorial}.
Every member of $A$ preserves the peripheral conjugacy
classes, as does its inverse. Applying the actual
automorphism and the individual conjugations in (B.36.2)
therefore preserves this exclusion. The limiting tuples
are pro-simple; they need not be discrete simple curves.

Conjugacy in a profinite group is a closed relation:
it is the image of the compact map
$(h,x)\mapsto(x,hxh^{-1})$. Exponentiation
$(x,u)\mapsto x^u$ for $u\in\widehat{\mathbb Z}$ is
continuous. For a fixed $\gamma$, condition (B.36.1) is
therefore a finite union of closed conditions on a section
and a member of $\mathcal K$. Intersect over all $\gamma$.
The resulting incidence subset is closed, and projection
along the compact space $\mathcal K$ is closed.

This argument uses the usual section topology, with no
compactness assumption on the set of sections itself.
For completeness, view a section as its closed image
subgroup $H\subset E$, with the subgroup topology used in
\cite[Section 3]{B_CarlsonCorwin}. Evaluation at a fixed
$\gamma$ is continuous on these subgroups: for a convergent
net $H_j\to H$, every convergent subnet of the elements
$s_j(\gamma)$ has limit in $H$ lying above $\gamma$;
there is exactly one such element. Compactness of $E$
then gives $s_j(\gamma)\to s(\gamma)$. Quotienting by
the compact geometric conjugation group preserves the
closed conclusion. Conjugation changes the witness tuple
by an element of $A$, so the subset is invariant.

Suppose now that the restricted action is realized by a
smooth marked curve $U/L$, with $L/k$ finite. Centre-freeness
of $\Delta$ identifies its arithmetic extension with the
pullback
\[
 E_{s,L}=\operatorname{Aut}(\Delta)
       \mathop{\times}_{\operatorname{Out}(\Delta)}G_L
       \ \simeq\ \pi_1(U),
 \tag{B.36.3}
\]
where the second isomorphism is the specified peripheral
realization. Choose one group $I=\langle c_i\rangle$ from
the witness. Its conjugacy class has a finite orbit under
$G_L$. Its stabilizer $G_0$ is closed, since it is the image
of the compact normalizer $N_{E_{s,L}}(I)$, and is of finite
index, hence open. Within that normalizer put
\[
 N^{\pm\chi}=
 \{e:\operatorname{Ad}(e)(c_i)
            =c_i^{\pm\chi(p(e))}\},\qquad
 N^\chi=
 \{e:\operatorname{Ad}(e)(c_i)
            =c_i^{\chi(p(e))}\}.
 \tag{B.36.4}
\]
Condition (B.36.1), with its existential conjugator,
makes $N^{\pm\chi}\to G_0$ surjective. Both are closed
subgroups. The ratio of the conjugation character and
$\chi\circ p$ gives a continuous homomorphism
$N^{\pm\chi}\to\{1,-1\}$ with kernel $N^\chi$.
Consequently $p(N^\chi)$ has index at most two in $G_0$;
it is compact and thus open in $G_L$. This is an actual
closed subgroup of the arithmetic normalizer with open
arithmetic image acting exactly cyclotomically on $I$.
No splitting of this normalizer is assumed.

The cyclotomic-normalizer theorem, in the exact form
\ref{A:cyclotomic-normalizer}, or
\cite[Theorem 20]{B_StixCuspidal}, puts $I$ in a cusp
inertia group of $U$. Its number-field hypothesis is
verified by the incompatible Frobenius weights of a torus
and a Jacobian at a good finite prime. Under (B.36.3)
the cusp groups are the prescribed peripheral groups.
This contradicts the exclusion proved above, and proves
nonrealization over every finite extension.
\end{proof}

\begin{statement}{B.37}{Proved}{Stable-boundary cuspidal sections are non-admissible}
Let $k$ be a number field, $2g-2+n>0$, and set
$\mathcal M=\mathcal M_{g,n,k}$ and
$\overline{\mathcal M}=\overline{\mathcal M}_{g,n,k}$.
For a $k$-rational stable boundary object $x$, let
$\mathcal V_x$ be the intersection with $\mathcal M$ of
the henselian neighbourhood of its residual gerbe in
$\overline{\mathcal M}$. Define
\[
 C_{\mathcal M,\overline{\mathcal M}}
 =\bigcup_{x\in(\overline{\mathcal M}\setminus\mathcal M)(k)}
 \operatorname{Im}\bigl(
       \operatorname{Sec}(\mathcal V_x/k)
        \longrightarrow\operatorname{Sec}(\mathcal M/k)\bigr),
 \qquad
 C^{\mathrm{cl}}_{\mathcal M,\overline{\mathcal M}}
 =\overline{C_{\mathcal M,\overline{\mathcal M}}}.
 \tag{B.37.1}
\]
Isomorphic boundary objects give the same term. For schemes
this is the nearby-point definition of Carlson--Corwin;
for stacks the residual gerbe and its finite automorphisms
are retained as described below.

Every section in $C^{\mathrm{cl}}_{\mathcal M,\overline{\mathcal M}}$
has an outer action on the universal geometric punctured
curve group which is not realized by a smooth marked curve,
even after restriction to any finite extension of $k$.
In particular, these sections are non-admissible in the
constructive sense of scans 480 and 487--488, and none is
the section of an interior moduli point. The same conclusion
holds for unordered markings, with the associated finite
permutation quotient retained.
\end{statement}
\begin{proof}
We first specify the scheme and stack maps in (B.37.1).
For a smooth scheme pair $V\subset X$ with normal crossings
boundary and $x\in(X\setminus V)(k)$, Carlson--Corwin define
the packet using the actual morphism
\[
 X_{(x)}\times_X V\longrightarrow V.
 \tag{B.37.2}
\]
Their excision lemma makes this construction invariant under
\emph{\'etale} replacement with the same residue point.
Their local fundamental-group proposition identifies its
geometric kernel with the dual of the boundary-character
group, twisted by $\widehat{\mathbb Z}(1)$; their Theorem A
asserts independence of the \emph{closure} from the normal
crossings compactification \cite{B_CarlsonCorwin}.
We have read those precise statements and their local
comparison. The theorem is stated for schemes. For the
stack in this statement we keep its specified stable
compactification and now give the local stack construction.

The residual gerbe at the object $x$ is
$B_k\operatorname{Aut}(x)$. Its henselian neighbourhood
is formed by the inverse system of representable
\emph{\'etale} neighbourhoods inducing an isomorphism on
that gerbe. In characteristic zero its finite automorphism
group is \emph{\'etale} and linearly reductive. The local
quotient construction of
\cite[Theorem 3.2 and Proposition 3.6, Step 1]{B_AOV}
gives a finite \emph{\'etale} torsor
\[
 \operatorname{Spec}R\longrightarrow
 [\operatorname{Spec}R/A]
       =\overline{\mathcal M}_{(x)},
 \tag{B.37.3}
\]
in that henselian system. Here $R$ is regular henselian,
$A$ is the finite group scheme extending
$\operatorname{Aut}(x)$, and its closed fibre torsor
over $B_k\operatorname{Aut}(x)$ has total space
$\operatorname{Spec}k$. The corrected obstruction argument
in \cite{B_AOVCorrigendum} also gives (B.37.3); in the
present \emph{\'etale} characteristic-zero case no
nonreduced stabilizer occurs. Removing the pulled-back
boundary gives
\[
 V_R=\operatorname{Spec}R\setminus D_R
       \longrightarrow [V_R/A]=\mathcal V_x
       \longrightarrow\mathcal M.
 \tag{B.37.4}
\]
These are the actual maps used to define its packet.
Changing the presentation gives the same henselian
neighbourhood and the same section image. A finite
extension splitting $A$ and the boundary branches makes
the first map an ordinary finite-group torsor. A section
of its quotient lifts to that torsor after restriction to
an open subgroup of $G_k$, by taking the kernel of its
finite torsor monodromy. The finite quotient is not
discarded in (B.37.1).

Let $r>0$ be the number of nodes of the geometric stable
curve $x$. In an \emph{\'etale} deformation chart the
boundary smoothing parameters are $t_1,\ldots,t_r$;
near node $i$ the universal curve has equation
$uv=t_i$. After the finite extension which fixes the
nodes and their branches, the local base group has
geometric inertia
$T^r$, where $T=\widehat{\mathbb Z}(1)$.
The local node group has geometric inertia
$T^{r+1}$, and its map to the base is
\[
 (a,b,d_2,\ldots,d_r)
        \longmapsto (a+b,d_2,\ldots,d_r),
 \qquad
 I_i=\{(a,-a,0,\ldots,0):a\in T\}.
 \tag{B.37.5}
\]
The coordinate numbered $1$ in this display is $t_i$.
In particular the map is surjective and the kernel
is an injected copy of $T$.

Here injection into the full geometric curve group is
essential. It is exactly the node injection and addition
map of \cite[Lemma 2.11(ii)--(iii)]{F-HoshiMochizukiNodes},
with all primes allowed. The source's Definition 2.10
identifies the full geometric log fibre group with the
punctured smooth surface group and indexes the local
character monoid by all the nodes. The comparison with
the generic smooth curve and with its universal-family
outer action is the comparison of finite covering
categories: use \cite[Theorem 7.6]{F-Illusie} for the
interior and \cite[Corollary 3.9]{F-HoshiLog} for the
completed strictly local fibre. In characteristic zero
every ramification index here is tame. Passing from the
henselization to its completion preserves the residue
group and the normal-crossings character monoid, hence
induces an isomorphism on their local groups by the same
Kummer calculation. All these maps preserve the base
quotient and the node kernel (B.37.5).

For any section of the local base group, pull back the
surjective node-group map along that section. It gives
a closed subgroup of the corresponding arithmetic
curve group which projects onto $G_k$ and acts on
$I_i$ by $\chi$. Thus this assertion holds for every
parameter in the packet; it does not restrict that
parameter to the image of an actual tangent vector
or of a one-dimensional arc. Without the preliminary
field extension, arithmetic elements permute the nodes;
exchanging the two branches changes $a$ to $-a$ in
(B.37.5). An automorphism in the finite quotient $A$
has the same effect: it permutes nodes and branches,
and a branch-preserving local coordinate change
$u\mapsto u\cdot\text{unit}$ acts identically on its
inertia. Consequently the full local group, including
$A$, satisfies exactly the permutation and $\pm\chi$
formula (B.36.1).

Choose a complex embedding of $\bar k$ and a geometric
base point of $\mathcal M$. Analytically smoothing the
stable fibre identifies its node groups with a nonempty
multicurve in the punctured topological fibre. Stability
ensures that these curves are essential and
nonperipheral. Changing the path from this local fibre
to the chosen geometric base point acts through the
geometric monodromy group $Q$; changing the fibre base
paths gives the individual conjugations of (B.36.2).
The finite-cover comparisons just described identify
these as the same full profinite node groups. Therefore
each packet in (B.37.1) supplies a tuple belonging to
the compact space $\mathcal K$ of B.36, and its moduli
section lies in the closed obstruction subset proved
there. Its entire closure lies in that same subset.
This last compactness step prevents a node witness
from being lost when boundary packets vary.

Apply B.36 to $E=\pi_1(\mathcal M)$ and the outer
monodromy of the universal punctured curve. It excludes
a smooth peripheral realization after every finite
extension. A moduli point gives precisely such a
realization, by pulling back that universal curve.
The constructive admissible germs of the manuscript
are obtained from an actual tripod curve by finite
\emph{\'etale} covering and filling specified holes;
those operations produce smooth marked curves and
their actual outer actions. Thus a boundary section
cannot be one of those germs. Passing to unordered
markings only adds the finite permutations already
allowed in the incidence relation and in (B.37.4), so
the same proof applies.

This is a necessary closed obstruction. It does not
assert that every section outside the closed boundary
set is a moduli point, or that the absence of these
pro-simple cyclotomic witnesses is a sufficient
realization criterion.
\end{proof}

\begin{statement}{B.38}{Proved using the stipulated local theorem}{Unramification over a p-adic base curve}
Let $k/\mathbb Q_p$ be finite, let $B/k$ be a smooth
geometrically connected hyperbolic curve, put $K=k(B)$,
and let $f:X\to B$ be a smooth proper family of
geometrically connected curves of genus at least two.
Write $E_X=\pi_1(X)$, $E_B=\pi_1(B)$, and let $\Delta$
be the geometric generic fibre group. The actual
generic-fibre comparison is
\[
 \pi_1(X_K)\ \xrightarrow{\ \sim\ }\
 E_X\times_{E_B}G_K.
 \tag{B.38.1}
\]
Pullback along $G_K\twoheadrightarrow E_B$ induces a
bijection
\[
 \operatorname{Sec}(E_X\to E_B)
 \ \xrightarrow{\ \sim\ }\
 \operatorname{Sec}(\pi_1(X_K)\to G_K),
 \tag{B.38.2}
\]
where both sides are taken modulo conjugation by the
same group $\Delta$. Every section on the left is
point-theoretic over $k$: for each finite extension
$l/k$ and each $b\in B(l)$, its restriction to the
point section at $b$ is induced by a unique point of
$X_b(l)$.

The maps on algebraic sections and on fundamental-group
sections fit in the commutative diagram
\[
\begin{array}{ccc}
 \operatorname{Sect}_B(X)&\xrightarrow{\sim}&X_K(K)\\
 \downarrow&&\downarrow\\
 \operatorname{Sec}(E_X\to E_B)&\xrightarrow{\sim}&
       \operatorname{Sec}(\pi_1(X_K)\to G_K).
\end{array}
\tag{B.38.3}
\]
The upper bijection is extension by properness. This
statement does not assert surjectivity of either
vertical map for a nonconstant family.
\end{statement}
\begin{proof}
The homotopy sequence for a relative hyperbolic curve
over this base is short exact:
\[
 1\longrightarrow\Delta\longrightarrow E_X
       \longrightarrow E_B\longrightarrow1.
 \tag{B.38.4}
\]
One can obtain its left exactness from the homotopy
fibre comparison and the fact that $B$ is an \'etale
$K(\pi,1)$. More generally, the kernel of the fibre
map lies in the centre of the hyperbolic fibre group
and is therefore trivial. These are precisely the
quasifibration comparison and the centre-free argument
in the proof comparing total and generic-fibre
$\pi_1$-openness in the supplied primary source
\cite{F-HolzschuhSchmidtStix}. The proper smooth fibre
comparison identifies $\Delta$ with the full geometric
group of every fibre, not merely its abelianization.
Apply the same exact sequence over $K$. The natural
map from $\pi_1(X_K)$ to the pullback in (B.38.1)
is the identity on $\Delta$ and on the quotient
$G_K$. The two exact sequences prove that it is an
isomorphism. This also identifies the two conjugating
groups in (B.38.2).

Let $s:G_K\to\pi_1(X_K)$ be a section and let
$h:G_K\to E_X$ be its first component under (B.38.1).
We prove that $h$ kills every inertia group belonging
to a closed point of $B$. Fix such a point $b$, with
residue field $k_b$. The completed local field is
$K_b\simeq k_b((t))$; choose a coefficient field and
a uniformizer. Its decomposition map to $E_B$
factors through the actual point section
$G_{k_b}\to E_B$, because a finite \'etale cover of
$B$ is unramified at $b$. In the proper smooth family,
the actual preimage of this point-section subgroup is
\[
 E_X\times_{E_B}G_{k_b}\ \simeq\ \pi_1(X_b).
 \tag{B.38.5}
\]
Indeed the map from the fibre group gives the identity
on the proper smooth geometric fibre group and on
$G_{k_b}$ in the two exact sequences. In particular,
the restriction of $h$ gives a homomorphism
\[
 G_{k_b((t))}\longrightarrow\pi_1(X_b)
                 \quad\text{over }G_{k_b}.
 \tag{B.38.6}
\]

The valuation in (B.38.6) has residue characteristic
zero. Its entire inertia is tame, and the chosen
coefficient field and compatible roots of $t$ give
\[
 G_{k_b((t))}\simeq
 \widehat{\mathbb Z}(1)\rtimes G_{k_b}
 \simeq\pi_1(\mathbb G_{m,k_b}),
 \tag{B.38.7}
\]
with the standard section of $\mathbb G_m$ at $1$.
Now precompose (B.38.6) with the map induced by the
actual open immersion
\[
 \mathbb P^1_{k_b}\setminus\{0,1,\infty\}
           \longrightarrow\mathbb G_{m,k_b}.
 \tag{B.38.8}
\]
Its map on geometric groups surjects onto
$\widehat{\mathbb Z}(1)$. The image of that group
in the proper hyperbolic fibre group is therefore
procyclic. The finite-$p$-adic-field version of
Proposition~\ref{F-local-dichotomy}, whose proof
uses the stipulated local section theorem, says
that this image is open, trivial, or contained in a
nontrivial cusp inertia. Properness eliminates the
last alternative. The open alternative is impossible:
an open subgroup of the genus-at-least-two proper
geometric group is the fundamental group of a finite
\'etale cover of genus at least two, and its
abelianization has $\mathbb Z_\ell$-rank at least
four for every prime $\ell$. A procyclic group has
rank at most one. Thus (B.38.6) kills the entire
$t$-adic inertia.

The kernel of $G_K\twoheadrightarrow E_B$ is the
closed normal subgroup generated by these inertia
groups. This follows directly from the finite-cover
description: a finite separable extension unramified
at every closed point of the regular curve $B$ has
finite \'etale normalization over $B$, and every
finite \'etale cover has that unramified property.
Consequently $h$ factors uniquely through a continuous
map $\bar h:E_B\to E_X$. The equation $f_*h=p_B$
and surjectivity of $p_B:G_K\to E_B$ make $\bar h$
a section. Conversely, pull back any such $\bar h$
using (B.38.1). These constructions are inverse on
actual homomorphisms and commute with conjugation by
$\Delta$, proving (B.38.2). Keeping (B.38.1) is
essential here: it proves that the pulled-back section
is the original $s$, rather than only a section with
the same image in $E_X$.

After any finite extension $l/k$, restrict a relative
section to the actual point section of $b\in B(l)$.
The same fibre comparison gives a section of
$\pi_1(X_b/l)$. The stipulated proper local theorem
gives its unique realizing point in $X_b(l)$. Thus
the map $E_B\to E_X$ has the asserted point-theoretic
property over every finite extension, with fibre
membership guaranteed by (B.38.5).

Finally, a $K$-point of the proper generic fibre
extends over every local discrete valuation ring of
$B$ by properness; uniqueness by separatedness glues
these extensions to a section of $X\to B$.
Restriction is its inverse, and functoriality gives
(B.38.3). Every smooth proper hyperbolic curve over a
one-variable function field over a finite $p$-adic
field admits such a model after replacing the base
by a nonempty hyperbolic open. The missing general
step is therefore the realization of these relative,
point-theoretic sections for nonconstant families.
The dominant relative Hom theorem in
\cite{F-HolzschuhSchmidtStix} does not supply that
step: the generic map from a point to a curve is not
dominant, and the corresponding section subgroup is
not open. A map on fundamental groups also does not
by itself provide a map of the full \'etale topoi.
No such extension of its hypotheses is used here.
\end{proof}

\begin{statement}{B.39}{Proved using the stipulated local theorem}{Geometric-image trichotomy over every sub-p-adic field}
Let $K$ be a sub-$p$-adic field, let $V/K$ be a smooth
geometrically connected variety, let $U/K$ be a
hyperbolic curve, and let
$h:\pi_1(V)\to\pi_1(U)$ be a continuous
$G_K$-homomorphism. Write $P_V,P_U$ for the geometric
groups. Exactly one of the following holds:
\[
 h(P_V)\text{ is open in }P_U;\qquad
 h(P_V)=1;\qquad
 1\ne h(P_V)\subset I_c
 \text{ for a unique lifted cusp }c.
 \tag{B.39.1}
\]
In the third case $h(\pi_1(V))\subset D_c$, and the
underlying cusp is $K$-rational. In the second case
$h$ factors through an actual section
$G_K\to\pi_1(U)$; no point realization of that section
is asserted.

The transfer used here preserves the full geometric
homomorphism. More precisely, there exist an extension
$K\hookrightarrow F$, a countable field $F_0\subset F$,
an embedding $F_0\hookrightarrow l$ into a finite
extension of $\mathbb Q_p$, and a descended datum
$(V_0,U_0,h_0)/F_0$ whose base change to $F$ is
$(V_F,U_F,h_F)$. The base change to $l$ has the same
geometric homomorphism through the characteristic-zero
base-change isomorphisms, including the peripheral
subgroups of the target. These are the actual field
maps used; an embedding of $F$ into $l$ is unnecessary.
\end{statement}
\begin{proof}
By definition of sub-$p$-adic, choose $F$ finitely
generated over $\mathbb Q_p$ containing $K$. Base
change gives
\[
 \pi_1(V_F)=\pi_1(V)\times_{G_K}G_F,
 \qquad
 \pi_1(U_F)=\pi_1(U)\times_{G_K}G_F.
 \tag{B.39.2}
\]
Here and below the compatible geometric base points
are understood. The geometric comparison maps are
isomorphisms in characteristic zero; one precise
primary reference is \cite[Theorem 1.1 and its
characteristic-zero specialization]{B_LandesmanBaseChange}.
Thus (B.39.2) preserves the geometric homomorphism,
without requiring $G_F\to G_K$ to be injective or open.

We give the countable descent of $h_F$ explicitly.
Put $E=\pi_1(U_F)$ and $\Delta=P_{U_F}$, and choose
a decreasing cofinal sequence of characteristic open
subgroups $N_i\subset\Delta$. Such a sequence exists
because $\Delta$ is finitely generated profinite:
intersect its finitely many subgroups of index at
most $i$. Each $N_i$ is normal in $E$.
In $E/N_i$ choose an open normal subgroup $J_i$
intersecting its finite subgroup $\Delta/N_i$
trivially. At each next stage choose $J_{i+1}$ inside
the inverse image of $J_i$. Set
\[
 M_i=(E/N_i)/J_i,\qquad
 Q_i=G_F/p(J_i),
 \qquad F_i=\overline F^{\,p(J_i)}.
 \tag{B.39.3}
\]
The finite Galois group of $F_i/F$ is $Q_i$.
Restriction of $p$ identifies $J_i$ with $p(J_i)$,
and the quotient maps give an isomorphism of
extensions
\[
 E/N_i\ \simeq\ G_F\times_{Q_i}M_i.
 \tag{B.39.4}
\]
For surjectivity, lift a prescribed element of $M_i$
to $E/N_i$ and correct its arithmetic projection by
the unique element of $J_i$ over the required
element of $p(J_i)$. Injectivity follows from
$J_i\cap(\Delta/N_i)=1$.

Equivalently, (B.39.4) is the arithmetic group of
the finite \'etale gerbe
\[
 \mathcal G_i=[\operatorname{Spec}F_i/M_i],
 \tag{B.39.5}
\]
where $M_i$ acts through $Q_i$. The quotient of $E$
and its composite with $h_F$ give actual maps
$U_F\to\mathcal G_i$ and $V_F\to\mathcal G_i$.
Concretely, these are finite $M_i$-torsors on the
two varieties, with equivariant maps to
$\operatorname{Spec}F_i$. The compatibility of
$J_i$ gives transition maps of these torsors and
gerbes. The inverse limit of the gerbes is
$\Pi^{\mathrm{\acute et}}_{U_F/F}$: it has precisely
the cofinal geometric kernels $N_i$, and its base
quotient remains $G_F$ at every stage. The exact
finite-gerbe and fundamental-group comparison is
also described in \cite[\S8 and Lemma 9.4]{B_BorneVistoli}.

Every item just listed has finite presentation:
the two varieties, the smooth compactification and
boundary of $U_F$, the finite \'etale algebras $F_i$,
the finite torsors, their group actions, transition
maps, and compatibility isomorphisms. Use finite
affine covers to describe them by equations and
maps. There are countably many items and equations.
They therefore descend together to a countable
subfield $F_0\subset F$. The gerbe formulation
does not omit the coherence data: their finitely
presented groupoid torsors and the corresponding
isomorphisms are included. The finite-presentation
limit principle used in this description is
\cite[Lemma 102.3.7, Tag 0CMX]{B_StacksLimits};
its groupoid-descent proof applies to the objects,
arrows, and their equations. All torsor properties
and equations hold after the faithfully flat
extension $F/F_0$, and hence hold over $F_0$.

We can impose a useful form on this countable
field. Let $k$ be the relative algebraic closure
of $\mathbb Q_p$ in $F$; it is a finite extension
of $\mathbb Q_p$. Choose a separating transcendence
basis $t_1,\ldots,t_d$ and one primitive element
$z$ for the finite extension $F/k(t_1,\ldots,t_d)$.
Write every coefficient of the countable data as
a rational expression in $t_1,\ldots,t_d,z$ with
coefficients in $k$. Include also the coefficients
of the minimal polynomial of $z$. The subfield
$k_0\subset k$ generated by these countably many
constants is countable, and we may take
\[
 F_0=k_0(t_1,\ldots,t_d,z)\subset F.
 \tag{B.39.6}
\]
The polynomial defining $z$ remains irreducible
over $k_0(t_1,\ldots,t_d)$ because it is irreducible
over $k(t_1,\ldots,t_d)$. Thus (B.39.6) preserves
all the algebraic relations of the descended data.

Let $U_0,V_0,\mathcal G_{i,0}$ denote these descents.
Their compatible maps give
\[
 V_0\longrightarrow\varprojlim_i\mathcal G_{i,0},
 \qquad
 \Pi^{\mathrm{\acute et}}_{U_0/F_0}
       \longrightarrow\varprojlim_i\mathcal G_{i,0}.
 \tag{B.39.7}
\]
The second map is an isomorphism. To check this,
pass to an algebraic closure of $F_0$ embedded in
an algebraic closure of $F$. Characteristic-zero
invariance identifies the geometric group of
$U_0$ with $\Delta$, and the kernels of its maps
to the geometric groups of $\mathcal G_{i,0}$
are precisely the cofinal $N_i$. Consequently
(B.39.7) is an isomorphism on geometric groups.
The arithmetic quotient is the identity of
$G_{F_0}$, so the exact fundamental sequences
give the asserted isomorphism. The first map
therefore defines an actual $h_0$ over $G_{F_0}$,
whose base change to $F$ is $h_F$, up to the usual
geometric base-path conjugacy. This construction
produces one compatible homomorphism; it does not
replace it by unrelated homomorphisms at separate
finite levels.

We next embed the particular field (B.39.6) into
one finite $p$-adic field. After inverting a
nonzero polynomial and the discriminant, it is
the fraction field of an integral algebra $A$
finite \'etale over
$k_0[t_1,\ldots,t_d,1/g]$. The same presentation
over $k$ has function field $F$. Choose a closed
point of $\operatorname{Spec}(A\otimes_{k_0}k)$;
its residue field $l$ is finite over $k$. The
\'etale map is a $p$-adic analytic isomorphism on
sufficiently small neighborhoods of that point
and its image. In the resulting closed ball
in $l^d$, choose coordinates algebraically
independent over $k_0$. This is possible by the
Baire theorem: there are countably many nonzero
polynomials over the countable field $k_0$, and
each has a closed zero set with empty interior.
The lifted point defines $A\to l$ whose kernel
contracts to zero in $k_0[t_1,\ldots,t_d,1/g]$.
Since $A$ is integral and finite over that ring,
the kernel is zero: after passage to its fraction
field the finite integral domain is a field.
Thus evaluation extends to the field embedding
$F_0\hookrightarrow l$. For $d=0$ the same
argument is simply a closed point of the finite
\'etale algebra. In particular, the embedding
preserves all the countably many retained
relations simultaneously.

Base change of $h_0$ to $l$ has exactly the original
geometric homomorphism, through (B.39.2) and
(B.39.7). The descended compactification and
boundary identify the peripheral subgroups as
well. Apply the finite-local-field trichotomy
proved in Proposition~\ref{F-local-dichotomy}.
Its proof uses the stipulated local section
theorem, with the provenance recorded in B.7.
This proves the alternatives in (B.39.1).
In the nontrivial peripheral case, return to
the original group over $K$: its geometric
image $H$ is normal in $h(\pi_1(V))$, and the
full peripheral normalizer identity gives
\[
 h(\pi_1(V))\subset
 N_{\pi_1(U)}(H)=D_c.
 \tag{B.39.8}
\]
The arithmetic projection of this image is all
of $G_K$, so the underlying cusp is rational.
If $H=1$, the quotient $\pi_1(V)\twoheadrightarrow
G_K$ gives the asserted section factorization.
Finally, a nontrivial cusp inertia has infinite
index in the hyperbolic geometric group, so
the three alternatives are disjoint.
\end{proof}

\begin{statement}{B.40}{Proved using the stipulated local theorem}{Unramification over every smooth p-adic model}
Let $k/\mathbb Q_p$ be finite, let $S/k$ be a smooth
geometrically connected variety, put $K=k(S)$, and
let $X\to S$ be a smooth proper family of
geometrically connected curves of genus at least two.
With $\Delta=\pi_1(X_{\overline K})$, the actual
generic-fibre comparison and restriction map are
\[
 \pi_1(X_K)\simeq
       \pi_1(X)\times_{\pi_1(S)}G_K,
 \qquad
 \operatorname{Sec}(\pi_1(X)\to\pi_1(S))
       \xrightarrow{\sim}
 \operatorname{Sec}(\pi_1(X_K)\to G_K).
 \tag{B.40.1}
\]
Both section sets use conjugation by $\Delta$.
Every relative section on the left restricts at
each $s\in S(l)$, for every finite $l/k$, to the
section of a unique point in $X_s(l)$.
The statement makes no assertion that every such
relative section is induced by a morphism $S\to X$.
\end{statement}
\begin{proof}
The centre-free fibre argument used in B.38 gives
\[
 1\longrightarrow\Delta\longrightarrow\pi_1(X)
       \longrightarrow\pi_1(S)\longrightarrow1
 \tag{B.40.2}
\]
over this arbitrary-dimensional normal base; no
$K(\pi,1)$ assumption on $S$ is required for the
centre-free argument. Comparing with the exact
sequence over $K$ proves the first isomorphism in
(B.40.1), with the identity on its kernel and quotient.
Let $h:G_K\to\pi_1(X)$ be the first component of
a generic section.

For a codimension-one point $\eta_D$ of $S$, let
$F_D=k(\eta_D)$. This is a finitely generated
extension of $\mathbb Q_p$. Completion of its
regular local discrete valuation ring gives an
equal-characteristic-zero field
$K_D\simeq F_D((t))$. Its map to $\pi_1(S)$
factors through the residue map $G_{F_D}\to\pi_1(S)$.
The exact comparison needed here is the fibre
product
\[
 \pi_1(X_D)\xrightarrow{\sim}
       \pi_1(X)\times_{\pi_1(S)}G_{F_D},
 \qquad X_D=X\times_S\operatorname{Spec}F_D.
 \tag{B.40.3}
\]
It follows by proper smooth identification of the
geometric fibre groups and the two exact sequences
over $G_{F_D}$. In particular it does not require
$G_{F_D}\to\pi_1(S)$ to be injective.
Pairing the restriction of $h$ with the residue
projection gives the actual homomorphism
\[
 G_{F_D((t))}\longrightarrow\pi_1(X_D)
                                      \quad/G_{F_D}.
 \tag{B.40.4}
\]
Identify its source with $\pi_1(\mathbb G_{m,F_D})$
as in (B.38.7), and precompose the tripod-to-$\mathbb G_m$
map. B.39 applies over $F_D$. The geometric image
is procyclic, whereas every open subgroup of the
proper hyperbolic fibre group has abelianization
of rank at least four. Properness supplies no cusp
alternative. Thus (B.40.4) kills all its inertia.

By purity of branch locus, the kernel of
$G_K\twoheadrightarrow\pi_1(S)$ is the closed
normal subgroup generated by these codimension-one
inertias. Indeed, the normalization of the regular
finite-type scheme $S$ in a finite separable
extension is finite, and being unramified in
codimension one makes it finite \'etale by
\cite[Lemma 58.21.4 and Example 58.13.3]{B_StacksPurity}.
Every finite \'etale cover has the converse
unramified property. Therefore $h$ factors uniquely
through $\pi_1(S)$. The equation over $\pi_1(S)$
makes this factor a section, and the first
isomorphism in (B.40.1) recovers the original
generic section. This proves the second isomorphism,
including its conjugacy quotient.

Finally, the same exact fibre product for a point
$s:\operatorname{Spec}l\to S$ identifies the
restricted section with a section of the proper
hyperbolic curve $X_s/l$. The stipulated local
theorem gives its unique point in $X_s(l)$.
Every smooth proper hyperbolic curve over a
finitely generated extension of $\mathbb Q_p$
spreads over a nonempty smooth model of the above
kind. The remaining generic point-realization
problem is retained; the factorization and its
fibrewise point data alone do not construct that
point or a full \'etale-topos morphism.
\end{proof}

\begin{statement}{B.41}{Proved using the stipulated local theorem}{The complete isotrivial p-adic function-field case}
Let $K$ be finitely generated over $\mathbb Q_p$,
and let $U/K$ be a hyperbolic curve which becomes
constant after a finite extension: there exist
$K'/K$ finite, a finite $p$-adic field $k'\subset K'$,
a hyperbolic curve $U_0/k'$, and an isomorphism
\[
 U_{K'}\simeq U_0\times_{k'}K'.
 \tag{B.41.1}
\]
For affine $U$ this hypothesis concerns the whole
open curve, including its boundary. Let
$\overline U$ be its smooth compactification and
$D=\overline U\setminus U$. Then the natural map
\[
 U(K)\ \sqcup\!\coprod_{c\in D(K)}
       \operatorname{Sec}(D_c/G_K)/I_c
       \ \xrightarrow{\ \sim\ }\
       \operatorname{Sec}(\pi_1(U)\to G_K)
 \tag{B.41.2}
\]
is bijective. Each packet is the usual torsor under
$H^1(K,\widehat{\mathbb Z}(1))=
\varprojlim_n K^\times/K^{\times n}$.
In particular every proper isotrivial hyperbolic
curve over such a field satisfies point-section
surjectivity in every transcendence degree.
\end{statement}
\begin{proof}
First let $C=C_0\times_k K$ be proper, with
$C_0/k$ of genus at least two and $k$ finite over
$\mathbb Q_p$. Enlarge $k$ to its relative
algebraic closure in $K$ and base change $C_0$;
then choose a smooth geometrically connected
$k$-model $S$ of $K$. Apply B.40 to the constant
family $S\times_k C_0$. Its arithmetic product
formula identifies a relative section with an
actual homomorphism
\[
 a:\pi_1(S)\longrightarrow\pi_1(C_0)
                                          \quad/G_k.
 \tag{B.41.3}
\]
The finite-local-field trichotomy of
Proposition~\ref{F-local-dichotomy} makes its
geometric image open or trivial, since $C_0$ is
proper. In the open case Mochizuki's Hom theorem
gives a unique dominant $k$-morphism
$S\to C_0$. Its source may have any dimension:
this is Theorem 16.5 for a smooth pro-variety,
specialized to a smooth variety, followed by the
full-profinite Remark \cite{F-Mochizuki}.
In the trivial case (B.41.3) factors through
$G_k$, and the stipulated local theorem gives
the unique constant point of $C_0(k)$ inducing
that factor. In both cases the generic value of
the resulting morphism is a point of $C(K)$.
The cartesian comparison in B.40 proves that
its section is the original generic section.

Uniqueness can be checked independently on the
same model. Two $K$-points spread to morphisms
from a common nonempty open of $S$ to $C_0$.
Equality of their generic section classes gives
equality of the induced homomorphism classes on
that model, because $G_K$ surjects onto its
fundamental group. Restriction at each closed
point gives equal point sections over a finite
extension of $k$, hence equal points by local
injectivity. Closed points are dense in this
finite-type $k$-scheme, and the equality locus
of the two morphisms is closed because $C_0$ is
separated. The morphisms, and their generic
points, are therefore equal.

We also record the geometric centralizer needed
for finite descent. For a section of this constant
proper curve, a geometric element centralizing
its image centralizes the image of (B.41.3):
the map from $G_K$ to $\pi_1(S)$ is surjective,
and the map to $\pi_1(C_0)$ is the identity on
the geometric fibre group. When (B.41.3) has
open geometric image, its geometric centralizer
is trivial by slimness of the hyperbolic curve
group. When it has trivial geometric image,
it is the section of a $k$-point and has trivial
geometric centralizer by B.10 over the finite
$p$-adic field $k$. Thus every section of the
constant proper curve over $K$ has trivial
geometric centralizer.

Now let $C/K$ be proper and satisfy (B.41.1).
Replace $K'$ by a finite Galois extension over
$K$. The result just proved gives a unique point
$x'\in C(K')$ inducing the restriction of any
given section $s$ over $K$. Its uniqueness implies
that every element of $\operatorname{Gal}(K'/K)$
fixes $x'$, since the restricted section class is
invariant under that action. Thus $x'$ descends
to $x\in C(K)$. There is still a section descent
step to check. Conjugate $s$ so that it agrees
with $s_x$ on $G_{K'}$. For $g\in G_K$, the
element $s(g)s_x(g)^{-1}$ lies in the geometric
group and centralizes $s_x(G_{K'})$, by normality
of $G_{K'}$ and equality of the restrictions.
The centralizer vanishes by the constant case
over $K'$, so $s=s_x$. This proves the proper
classification and its uniqueness for every
isotrivial curve and every finite extension of
the original field.

It remains to pass to affine $U$, allowing the
compactification to have genus zero or one.
Every geometrically connected finite \'etale
cover of an isotrivial $U$ has an isotrivial
smooth compactification. To see the precise
descent, make (B.41.1), embed an algebraic closure
of $k'$ in an algebraic closure of $K'$, and
use the equivalence of finite \'etale covering
categories given by characteristic-zero base
change \cite{B_LandesmanBaseChange}. The
geometric covering descends to $\overline{k'}$,
and then to a finite extension of $k'$ because
the cover has finite presentation. Its isomorphism
with the original cover descends to a finite
extension of $K'$. Normalization of the smooth
compactification commutes with this isomorphism.
This proves exactly the needed isotriviality of
each compactified neighborhood.

For a section of $U$, take the cofinal
neighborhoods with compact genus at least two
constructed in B.7. That construction is
geometric: a high-genus geometric cover,
a characteristic subgroup contained in its
group, and the arithmetic neighborhood
$N s(G_K)$. In each such compactification the
proper isotrivial result just proved supplies
a unique $K$-point. Under a map of neighborhoods
the points agree, because they induce the same
proper section and the proper point map is
injective. Thus they form one compatible point
in the compactified neighborhood tower. Its
image in $\overline U$ belongs to $\overline U(K)$.
An interior image gives the point section of
$U$; a boundary image gives a section in its
actual cusp decomposition group, by the same
stabilizer calculation as B.7. This proves
surjectivity in (B.41.2).

For injection and disjointness, every proper
curve used in the neighborhood proof of
\cite[Theorem 14]{B_StixCuspidal} is itself the
compactification of a finite \'etale cover of
$U$. It is isotrivial by the comparison above.
The proper injectivity just proved therefore
applies at each such level and after every
finite extension of $K$. Explicitly,
equality of sections identifies the compact
centres in every neighborhood, and a conjugator
between two sections of one packet fixes the
lifted cusp at every finite level, hence lies
in its inertia. This is the actual map from
inertia-conjugacy to geometric-conjugacy classes
in (B.41.2). Its torsor is the Kummer torsor
of B.3 over $K$, with the displayed power
completion. No compactness of $K$ or of the
finite-level multiplicative quotients is used.
\end{proof}

\begin{statement}{B.42}{Proved}{A compatible tower invisible on all closed fibres}
Let $k$ be a finite extension of $\mathbb Q_p$,
choose $q\in k$ with $0<|q|<1$, and let $E_q/k$
be the Tate elliptic curve. Put $B=E_q\setminus\{O\}$.
For every integer $n>1$ there is a geometrically
connected finite \'etale cyclic cover
\[
 f_n:C_n\longrightarrow B,\qquad \deg(f_n)=n,
 \tag{B.42.1}
\]
whose fibre over every closed point $b\in B$ is
the disjoint union of $n$ copies of
$\operatorname{Spec}k(b)$. Nevertheless, its
generic fibre has no $k(B)$-point, and $f_n$
has no section over any nonempty open of $B$.
The covers are compatible under divisibility
and determine a continuous surjection
\[
 \omega:\pi_1(B)\longrightarrow\widehat{\mathbb Z}
 \tag{B.42.2}
\]
which is surjective on $\pi_1(B_{\bar k})$
and vanishes on every closed-point decomposition
group. Thus the map
\[
 \operatorname{Hom}_{\mathrm{cont}}
       (\pi_1(B),\widehat{\mathbb Z})
 \longrightarrow
 \prod_{b\in B^{\mathrm{cl}}}
 \operatorname{Hom}_{\mathrm{cont}}
       (G_{k(b)},\widehat{\mathbb Z})
 \tag{B.42.3}
\]
has a nonzero kernel.
\end{statement}
\begin{proof}
Tate's uniformization gives, functorially for
every finite extension $L/k$,
\[
 E_{q^n}(L)=L^\times/q^{n\mathbb Z},
 \qquad E_q(L)=L^\times/q^{\mathbb Z}.
\]
The identity on $L^\times$ induces the
$k$-isogeny
\[
 \bar f_n:E_{q^n}\longrightarrow E_q.
\]
This is the actual algebraic morphism
$\alpha_{n,1}$ in Tate's isogeny theorem:
take $t_1=q^n$, $t_2=q$, so that
$t_2^n=t_1$. The same theorem proves its
degree is $n$; its kernel is the constant
cyclic subgroup generated by the $k$-point
$[q]\in k^\times/q^{n\mathbb Z}$
\cite[Theorem 1 and the isogeny theorem]{B_TateUniformization}.
It is \'etale because $k$ has characteristic
zero. Define
\[
 C_n=E_{q^n}\setminus\bar f_n^{-1}(O),
 \qquad f_n=\bar f_n|_{C_n}.
\]
This is a geometrically connected open in
an elliptic curve, and (B.42.1) follows.

For a closed point $b$ put $L=k(b)$.
Represent its image in $E_q(L)$ by
$a\in L^\times$. Its complete inverse image
under $\bar f_n$ consists of
\[
 [a], [aq],\ldots,[aq^{n-1}]
             \quad\text{in }L^\times/q^{n\mathbb Z}.
\]
These are $n$ distinct $L$-points, all in
$C_n$ because $b\ne O$. This proves the
asserted equality of finite \'etale fibre
schemes, not merely nonemptiness.
On the other hand $C_n$ is integral and its
function field is a degree-$n$ extension
of $k(B)$. A point of its generic fibre would
give a $k(B)$-homomorphism from that extension
to $k(B)$. Such a homomorphism is injective
and is the identity on $k(B)$, which is
impossible for degree $n>1$. A section on
any nonempty open would give that generic
point, so no such section exists.

For $n\mid m$ the identity on the
uniformizing multiplicative group gives
$C_m\to C_n$, and translation by $q$ maps
to translation by $q$. Choose a point
$b_0\in B(k)$ represented by
$a_0\in k^\times\setminus q^{\mathbb Z}$,
and the compatible lifts represented by
this same $a_0$ on every $C_n$.
The deck quotients are then the compatible
maps $\pi_1(B)\to\mathbb Z/n\mathbb Z$.
Their restrictions to the geometric group
are onto because every $C_n$ is
geometrically connected. Their inverse
limit is (B.42.2); its image is compact
and maps onto every $\mathbb Z/n\mathbb Z$,
so both it and the geometric image are
all of $\widehat{\mathbb Z}$.
Every closed-point decomposition group
acts trivially on each fibre just computed,
so its image in every deck quotient,
and hence under $\omega$, is zero.
This proves (B.42.3). Tate also records
this same completely split infinite tower
explicitly \cite[p.~18 of the author manuscript]{B_TateUniformization}.
\end{proof}

\paragraph{Exact remaining realization step.}
The finite-cover local-point hypothesis in
\cite[Definition 0.2(vi)]{B_SaidiFunctionFields}
requires a generic point whenever all the
closed fibres of a finite \'etale curve cover
have residue-field points. B.42 disproves
that hypothesis for a finite $p$-adic base
field, even for a hyperbolic base curve.
Consequently it cannot be inserted into
the realization step after B.40.
The maps in B.42 are quotient maps of
fundamental groups; they are not sections
of a proper hyperbolic family. No
nongeometric section of such a family is
constructed by this example.
The unresolved assertion is still whether
the particular relative sections produced
in B.40 all arise from algebraic sections
for nonisotrivial families. B.41 resolves
that assertion in the stated isotrivial
range only.
